%================================================================
\documentclass[12pt]{amsart}
\newcommand{\deleted}[1]{}
\newcommand{\delete}[1]{}
\newcommand{\mynotes}[1]{}
\newcommand\notes[1]{}
\usepackage{txfonts}
\usepackage{amssymb}
\usepackage{eucal}
\usepackage{graphicx}
\usepackage{amsmath}
\usepackage{amscd}
\usepackage[all]{xy}
\usepackage[bookmarksnumbered=true,
            bookmarksopen=true,
            colorlinks,
            pdfborder=001,
            linkcolor=blue,
            anchorcolor=blue,
            citecolor=blue
            ]{hyperref}

\usepackage{amsthm}
\usepackage{mathrsfs}
\usepackage{enumitem}
\usepackage{amsfonts,latexsym}
\usepackage{xspace}
%\usepackage{hyperref}
\usepackage{epsfig}
\usepackage{float}
\usepackage{color}
\usepackage{fancybox}
\usepackage{colordvi}
\usepackage{multicol}
\usepackage{tikz}
\usepackage{wasysym}
\usepackage{xspace}
\usepackage[active]{srcltx} %SRC Specials for DVI Searching
\usepackage{mathtools} %  for left up and down index?
\usepackage{tabularx}
\usepackage{stfloats}
\usepackage{makecell}
\usepackage{caption}
%\setlength{\lineskip}{0.15em}
%\setlength{\parskip}{0.5em}
%======================================================================
\renewcommand
\baselinestretch{1}    %was    1, 1.5 for double sp

\topmargin -.8cm
\textheight 22.8cm
\oddsidemargin 0cm
\evensidemargin -0cm
\textwidth 16.3cm

\newcommand\changed[1]{#1}
\changed{\newcommand\binomial[2]{{#1 \choose #2}}}

%%%%%%%%%%%%%%%%%%%%%%%% Statements
\newtheorem{theorem}{Theorem}[section]
\newtheorem{lemma}[theorem]{Lemma}
\newtheorem{obser}[theorem]{Observation}
\newtheorem{coro}[theorem]{Corollary}
\newtheorem{problem}[theorem]{Problem}
\newtheorem{fact}[theorem]{Fact}
\newtheorem{question}[theorem]{Question}
\newtheorem{conj}[theorem]{Conjecture}
\newtheorem{prop}[theorem]{Proposition}
\newtheorem{conject}[theorem]{Conjecture}
\theoremstyle{definition}
\newtheorem{defn}[theorem]{Definition}
\newtheorem{remark}[theorem]{Remark}
\newtheorem{exam}[theorem]{Example}
\newtheorem{claim}[theorem]{Claim}
\deleted{\newtheorem{theorem}{Theorem}[section]
\newtheorem{prop-def}{Proposition-Definition}[section]
\newtheorem{coro-def}{Corollary-Definition}[section]
\newtheorem{sublemma}[theorem]{Claim}
\newtheorem{propprop}{Proposed Proposition}[section]
}
%==========================================================================
\newcommand{\nc}{\newcommand}
\renewcommand{\frak}{\mathfrak}
\newcommand{\efootnote}[1]{}


%%%%%%% Redefine inequality signs
\renewcommand\geq{\geqslant}
\renewcommand\leq{\leqslant}
\renewcommand\preceq{\preccurlyeq}
\renewcommand\succeq{\succcurlyeq}
\renewcommand\bar[1]{\overline{#1}}
\newcommand{\lra}{\longrightarrow}
\newcommand{\llra}{\longleftarrow}
\newcommand{\lraf}[1]{\stackrel{#1}{\lra}}
\newcommand{\llaf}[1]{\stackrel{#1}{\llra}}
\newcommand{\raf}[1]{\stackrel{#1}{\ra}}
\newcommand{\laf}[1]{\stackrel{#1}{\lra}}
\newcommand{\ra}{\rightarrow}
\newcommand{\lan}{\langle}
\newcommand{\ran}{\rangle}

\newcommand{\dk}{{\rm dim_{_{k}}}}
\newcommand{\holim}{{\rm Holim}}
\newcommand{\hocolim}{{\rm Hocolim}}
\newcommand{\colim}{{\rm colim\, }}
\newcommand{\limt}{{\rm lim\, }}
\newcommand{\Add}{{\rm Add }}
\newcommand{\Prod}{{\rm Prod }}
\newcommand{\Tor}{{\rm Tor}}
\newcommand{\Cogen}{{\rm Cogen}}
\newcommand{\Tria}{{\rm Tria}}
\newcommand{\Loc}{{\rm Loc}}
\newcommand{\Coloc}{{\rm Coloc}}
\newcommand{\tria}{{\rm tria}}
\newcommand{\Con}{{\rm Con}}
\newcommand{\Sum}{{\rm Sum}}
\newcommand{\cpx}[1]{#1^{\bullet}}
\newcommand{\st}[1]{#1\mbox{{\rm -{\underline{mod}}}}}
\newcommand{\Da}[1]{{\mathscr D}^-(#1)}
\newcommand{\Db}[1]{{\mathscr D}^b(#1)}
\newcommand{\Kb}[1]{{\mathscr K}^b(#1\mbox{{\rm -proj}})}
\newcommand{\Ka}{{\mathscr K}^{-, b}(A\mbox{{\rm -proj}})}
\newcommand{\KL}{{\mathscr K}^{-, b}(\Lambda\mbox{{\rm -proj}})}
\newcommand{\Kbb}[1]{{\mathscr K}^{-, b}(\C)}
\newcommand{\K}{{\mathscr K}}
\newcommand{\Kc}{{\mathscr K}(\C)}
\newcommand{\Cc}{{\mathscr C}(\C)}
\newcommand{\Cb}[1]{{\mathscr C}^b(#1)}
\newcommand{\md}{\ensuremath{\mbox{{\rm -mod}}}}
\newcommand{\Modcat}{\ensuremath{\mbox{{\rm -Mod}}}}
\newcommand{\nf}[1]{$#1$-Nakayama-finite}
%=========================================================================
%%%%%%%%%%%%%%%%%%%% new symbols
\nc{\mrm}[1]{{\rm #1}}
\DeclareMathOperator{\Ker}{Ker}
\DeclareMathOperator{\Stmap}{\mathcal{S}}
\DeclareMathOperator{\Zmap}{\mathcal{Z}}
\DeclareMathOperator{\Klenj}{\K^{\ell<j}(A\proj)}

\nc{\Hom}{\mrm{Hom}}
\nc{\Ext}{\mrm{Ext}}
\nc{\End}{\mrm{End}}
\nc{\rad}{\mrm{rad}}
\nc{\Aut}{\mrm{Aut}}
\nc{\A}{\mathcal{A}}
\nc{\B}{\mathcal{B}}
\nc{\C}{\mathcal{C}}
\nc{\M}{\mathcal{M}}
\nc{\X}{X^{\bullet}}
\nc{\Y}{Y^{\bullet}}
\nc{\Z}{\mathbb{Z}}
\nc{\T}{T^{\bullet}}
\nc{\U}{U^{\bullet}}
\nc{\V}{V^{\bullet}}
\nc{\f}{f^{\bullet}}
\nc{\g}{g^{\bullet}}
\nc{\h}{h^{\bullet}}
\nc{\rr}{r^{\bullet}}
\nc{\s}{s^{\bullet}}
\nc{\Lv}{{\bf L}\nu}
\nc{\Lm}{{\bf L}\mu}
\nc{\Lx}{{\bf L}\xi}
\nc{\Pv}{{\bf P}\nu}
\nc{\w}{\widetilde}
\nc{\N}{\mathcal{N}}
\nc{\cw}[1]{\widetilde{#1}^{\bullet}}
\nc{\tr}{\mrm{tr}}
\nc{\pr}{\mrm{Pic}}
\nc{\Out}{\mrm{Out}}
\nc{\ch}{\mrm{char}}
\nc{\add}{\mrm{add}}
\nc{\RHom}{\mrm{{\bf R}Hom}}
\nc{\OT}{\otimes^{\bf L}}
\nc{\thick}{\mrm{thick}}
\nc{\cone}{\mrm{cone}}
\nc{\proj}{\mbox{{\rm -}proj}}
\nc{\Proj}{\mrm{Proj\,}}
\nc{\gl}{\mrm{gl}}
\nc{\im}{\mrm{im}\,}
\nc{\id}{\mrm{id}}
\nc{\ind}{\mrm{ind}}
\nc{\lcd}[1]{\mrm{lcd}(#1)}
\nc{\diag}{\mrm{diag}}
\nc{\pid}{-pseudo-identity\;}
\nc{\tp}[1]{\mrm{TrPic}(#1)}
\nc{\res}{\mrm{res}}
\nc{\PI}[1]{\mrm{PI}(#1)}
\nc{\uu}[1]{T(#1)}
\nc{\soc}{\mrm{soc}}
\nc{\va}{\varepsilon}
%=========================================================================

\begin{document}
%\renewcommand{\thefootnote}{\alph{footnote}}
\renewcommand{\thefootnote}{\alph{footnote}}
\setcounter{footnote}{-1}
\footnote{2020 Mathematics Subject Classification: 18G80, 13D09, 16E35, 16L60.}
\renewcommand{\thefootnote}{\alph{footnote}}
\setcounter{footnote}{-1}
\footnote{Keywords: Derived category, Derived equivalence, Pseudo-identity, Standard derived equivalence, Truncated triangle, Triangle functor.}

\title[Rickard's question on derived equivalences]{Rickard's question on standard derived equivalences}
%\title[Rickard's question on derived equivalence]{Not every derived equivalence is standard}
\author[Wei Hu]{Wei Hu}
\address{Wei Hu: School of Mathematical Sciences, Beijing Normal University, 100875 Beijing, China}
\email{huwei@bnu.edu.cn}
\author[Changchang Xi]{Changchang Xi}
\address{Changchang Xi: School of Mathematical Sciences, Capital Normal University, 100048 Beijing, P.R. China}
\email{xicc@cnu.edu.cn}
\author[Jin Zhang]{Jin Zhang}
\address{Jin Zhang: School of Mathematical Sciences, Laboratory of Mathematics and Complex Systems, MOE, Beijing Normal University, 100875 Beijing, China}
\email{zj\_10@lzu.edu.cn}

\begin{abstract}
In 1991 Rickard asked whether every derived equivalence of algebras over a common field is standard. We construct a series of examples of non-standard derived equivalences, and thus answer the question negatively. Moreover, we give a necessary and sufficient condition for the question to be true.
\end{abstract}

\maketitle
%========================================================================
%\hyphenpenalty=8000

\tableofcontents

\section{Introduction}
Let $A$ and $B$ be two finite-dimensional algebras over a field. Denote by $\Db{A}$ the bounded derived category of the category $A\md$ of finitely generated left $A$-modules. An equivalence $F: \Db{A}\ra \Db{B}$ is said to be a {\it derived equivalence} if $F$ is a triangle equivalence \cite{R0}. Moreover, if in addition $F$ is naturally equivalent to the derived tensor functor induced by a bounded complex of $B$-$A$-bimodules, then $F$ is said to be a {\it standard derived equivalence} \cite{R}. Clearly, standard derived equivalences form an important class of derived equivalences and have many significant applications, such as lifting or restricting to stable equivalences of Morita type and calculating homological invariants and groups (see \cite{R, RZ, Y}).

How can an arbitrary derived equivalence be related to a standard derived equivalence? Rickard showed in \cite[Corollary 3.5]{R} that, given  a derived equivalence $F: \Db{A}\ra \Db{B}$, there always exists a standard derived equivalence $\w{F}: \Db{A}\ra \Db{B}$ such that $F$ and $\w{F}$ agree on the full subcategory of $\Db{A}$ consisting of projective $A$-modules and $F(\X)\simeq \w{F}(\X)$ for every object $\X$ of $\Db{A}$. However, it remains  unclear whether $F$ and $\w{F}$ are naturally equivalent. Hence he proposed the following question \cite{R}.
\begin{question}{\rm \cite{R} }
 Are all derived equivalences standard?
\end{question}

We say that Rickard's question \emph{holds for} a finite-dimensional algebra $A$ over a field if every derived auto-equivalence of $\Db{A}$ is standard. It is known that Rickard's question holds for $A$ if and only if every derived equivalence $\Db{A}\ra \Db{B}$ between $A$ and an arbitrary algebra $B$ is standard.

Neeman \cite{N} investigated the question with the emphasis on maps between triangles and obtained a positive answer by adding axioms to the notion of triangulated categories. However, the question remains open in general. Only a few cases have been verified. For instance, a positive answer has been obtained for hereditary algebras \cite{MY}, triangular algebras \cite{C}, and derived-discrete algebras \cite{BC, CZ}.

In the present paper, we will provide a negative answer to Rickard's question. This is done by understanding triangle functors of special form. During this course, we are led to construct a series of counterexamples.

Following \cite{CY}, a derived auto-equivalence $(F, \xi): \Db{A}\ra \Db{A}$  (respectively,  $\Kb{A}\ra \Kb{A}$) is called a {\it pseudo-identity} if $F$ fixes all objects of $\Db{A}$ (respectively, $\Kb{A}$) and fixes $A\md[i]$ (respectively, $A\proj[i]$) for all $i\in \Z$.
%We write $\PI{A}$ for the group of all natural isomorphism classes of pseudo-identities $\Db{A}\ra \Db{A}$.

First, we give a criterion for Rickard's question to hold.

\begin{theorem}\label{t0}
Let $A$ be a finite-dimensional algebra over a field. Then Rickard's question holds for $A$ if and only if every pseudo-identity $(F, \id): \Kb{A}\ra \Kb{A}$ that fixes all truncated triangles is naturally isomorphic to the identity triangle functor.
\end{theorem}

For the definition of truncated triangles, we refer to Definition \ref{def-truncated}. It was proved in \cite[Theorem 6.1]{CY} that Rickard's question holds for $A$ if every pseudo-identity $(F, \xi): \Kb{A}\ra \Kb{A}$ is naturally isomorphic to the identity. By Theorem \ref{t0}, we do not have to check every pseudo-identity. It is enough to verify only those pseudo-identities  $(F, \id)$ that fix all truncated triangles. As an application of Theorem \ref{t0}, we give a short proof of Rickard's question for derived-discrete algebras in some cases; see Example \ref{ex}.

Based on Theorem \ref{t0} and a reduction lemma (see Lemma \ref{coro}), we have the following negative answer to Rickard's question.

\begin{theorem}\label{t}
There exist infinitely many non-standard derived equivalences between finite-dimensional algebras over fields.
\end{theorem}

The counterexamples are given in Example \ref{counterexa}. They occur over a field of characteristic $2$. We conjecture that Rickard's question has a positive answer for any field of characteristic $p\ne 2$.

\medskip
The paper is outlined as follows: In Section \ref{pre}, we give basic definitions and prove some preparatory results. In Sections \ref{aut} and \ref{example}, we prove Theorems \ref{t0} and \ref{t}, respectively. Also, we display explicit counterexamples in Section \ref{example}.


\section{Preliminaries}\label{pre}
In this section, we introduce basic notation and establish preparatory results for the proofs of the main results in the introduction.

Let $R$ be a ring with identity. Denote by $R\md$ the category of finitely generated left $R$-modules and $R\proj$ the full subcategory of $R\md$ consisting of projective $R$-modules.

Let $\C$ be an additive category. Denote by $\Cc$ the category of complexes over $\mathcal{C}$. Let $\Kc$ be the homotopy category of $\C$, and let $\K^b(\C)$ be the bounded homotopy category of $\C$. If $\C$ is a full subcategory of an abelian category $\A$, then $\Kbb{\C}$ denotes the homotopy category of bounded-above complexes over $\mathcal{C}$ with bounded homology. For two morphisms $\f, \g:\X \to \Y$ of complexes, we write $\f\sim \g$ if they are homotopic.

A morphism $f: X\ra Y$ in an additive category $\C$ is said to be {\it radical} if $\id_Y-fg$ is an isomorphism for any $g\in \Hom_{\C}(Y, X)$.  If $\C$ is a Krull-Schmidt category, then this is equivalent to saying that for every
indecomposable object $Z$ and every pair of morphisms
$h:Z\to X$ and $g:Y\to Z$, the morphism $gfh$ is not an isomorphism (see, for example, \cite[Corollary 2.10]{Krause}).

Assume that $\C$ is a Krull--Schmidt category. Then every morphism in $\C$ decomposes into an isomorphism and a radical morphism. More precisely, for every morphism $f: X\ra Y$ in $\C$, there are decompositions $X=X_1\oplus X_2$ and $Y=Y_1\oplus Y_2$ and automorphisms $s_X: X\ra X$ and $s_Y: Y\ra Y$ such that
$$s_Yfs_X^{-1}=\left(\begin{matrix}
s& 0\\
0&r
\end{matrix}\right): X_1\oplus X_2\lra Y_1\oplus Y_2$$
 where $s: X_1\ra Y_1$ is an isomorphism and $r:X_2\ra Y_2$ is a radical morphism.

We refer to \cite{H} for further definitions and basic facts about triangulated categories.

\subsection{Standard triangles and maps between standard triangles}\label{triangles}
Let $\C$ be an additive category. Let $\X$ and $\Y$ be complexes over $\C$. Suppose that
$$\f: \X=(X^i, d_X^i: X^i\ra X^{i+1})\lra \Y=(Y^i, d_Y^i:Y^i\ra Y^{i+1})$$
is a morphism of complexes. For each $i\in \Z$, we define a truncated complex $X^{\geq i}$ of $\X$ by
$$X^{\geq i}:\;\; \cdots \lra 0\lra X^i\lraf{d_X^i} X^{i+1}\lraf{d_X^{i+1}} X^{i+2}\lra \cdots $$
with $X^i$ in degree $i$. Similarly, we define $X^{\leq i}$, and $X^{[u, v]}=(X^{\geq u})^{\leq v}$ if $u\leq v$.
We can also define a morphism $f^{\geq i}: X^{\geq i}\ra Y^{\geq i}$ of complexes by $(f^{\geq i})^j=f^j$ for all $j\geq i$. Similarly, we define $f^{\leq i}: X^{\leq i}\ra Y^{\leq i}$ and $f^{[u, v]}: X^{[u, v]}\ra Y^{[u, v]}$ for $u\leq v$.

For simplicity, we use the same symbol $\f$ for its homotopy class in $\Kc$. There is a triangle
$$\Delta(\f): \X\stackrel{\f}\lra \Y\stackrel{\theta(\f)}\lra \cone(\f)\stackrel{\pi(\f)}\lra \X[1]$$
in $\Kc$, where
$$\cone(\f):=(X^{i+1}\oplus Y^i, \left(\begin{matrix}
-d_X^{i+1}& 0\\
f^{i+1}&d_Y^i
\end{matrix}\right): X^{i+1}\oplus Y^i\ra X^{i+2}\oplus Y^{i+1}), $$
$$ \theta(\f):=\left(\begin{smallmatrix}
0\\
1
\end{smallmatrix}\right) \text{\;and\;} \pi(\f):=(1, 0).$$
This triangle is called a \emph{standard triangle}. It is known that every triangle in $\Kc$ is isomorphic to a standard triangle. By a \emph{graded} morphism $\g: \X\ra \Y$ of complexes over $\C$, we mean a family of morphisms $g^i: X^i\ra Y^i$ in $\C$ indexed by $i\in \Z$. Note that a graded morphism need not be a morphism of complexes.



Given two morphisms $\f$ and $\g$ of complexes,  consider the following diagram
 \[(\divideontimes)\;\;\;\;\xymatrix@R=1.5em{
\X\ar[r]^{\f}&\Y\ar[r]^-{\theta(\f)}\ar[d]^{\cpx{r}}& \cone(\f)\ar[r]^{\pi(\f)}
\ar@{..>}[d]&\X[1]\ar[d]^{\cpx{s}[1]} \\
\U\ar[r]^{\g}&\V\ar[r]^-{\theta(\g)}& \cone(\g)\ar[r]^{\pi(\g)}&\U[1]
}\]
in $\Kc$ whose rows are standard triangles. In what follows, we write a graded morphism between $\cone(\f)$  and $\cone(\g)$ as $\left(\begin{smallmatrix}
\cpx{a}& \cpx{b}\\
\cpx{c}& \cpx{d}
\end{smallmatrix}\right)$, where its $i$-component is $\left(\begin{smallmatrix}
a^i& b^i\\
c^i& d^i
\end{smallmatrix}\right): X^{i+1}\oplus Y^i\ra U^{i+1}\oplus V^i$ with $a^i: X^{i+1}\to U^{i+1}, b^i:Y^i\to U^{i+1}, c^i:X^{i+1}\to V^i$ and $ d^i:Y^i\to V^i$ being morphisms in $\C$. The following lemma tells us how this morphism can be simplified up to homotopy when it makes one of the two squares in $(\divideontimes)$ commute. It also gives a criterion for the existence of a morphism $\cone(\f)\ra \cone(\g)$ in $\Kc$, making the squares in $(\divideontimes)$ commute.


\begin{lemma}\label{st}
Keep the notations above. Then the following statements hold.
\begin{enumerate}
  \item[$(1)$] If $\left(\begin{smallmatrix}
\cpx{a}& \cpx{b}\\
\cpx{c}& \cpx{d}
\end{smallmatrix}\right): \cone(\f)\ra \cone(\g)$ is a morphism of complexes such that one of the squares in $(\divideontimes)$ commutes, then it is homotopic to a morphism of the form $\left(\begin{smallmatrix}
\cpx{x}&0 \\
\cpx{y}&\cpx{z}
\end{smallmatrix}\right)$.
\item[$(2)$]  A graded morphism $\left(\begin{smallmatrix}
\cpx{a}& 0\\
\cpx{c}& \cpx{d}
\end{smallmatrix}\right): \cone(\f)\ra \cone(\g)$
is a morphism of complexes if and only if $\cpx{a}: \X[1]\ra \U[1]$ and $\cpx{d}: \Y\ra \V$ are morphisms of complexes such that $\cpx{d} \f\sim \g(\cpx{a}[-1])$ with $\cpx{c}[-1]$ as a homotopy.

\item[$(3)$] There is a morphism $\cpx{\phi}:\cone(\f)\ra \cone(\g)$ of complexes such that both squares in $(\divideontimes)$ commute if and only if there is a morphism $\h: \Y\ra \U$ such that $\g \h \f\sim \g\cpx{s}-\cpx{r}\f$. In this case, $\cpx{\phi}=\left(\begin{smallmatrix}
\cpx{s}[1]&0 \\
\cpx{c}[1]&\cpx{r}+\g\h
\end{smallmatrix}\right): \cone(\f)\ra \cone(\g)$ with $\cpx{c}$ a homotopy of $(\cpx{r}+\g\h)\f\sim \g\cpx{s}$.

\item[$(4)$]  If two morphisms $\cpx{\phi}, \cpx{\psi}: \cone(\f)\ra \cone(\g)$ of complexes make the squares in $(\divideontimes)$ commute, then there exists a morphism $\h: \Y\ra \U$ of complexes with $\g\h\f\sim 0$ such that
$\cpx{\phi}-\cpx{\psi}\sim\left(\begin{smallmatrix}
0&0 \\
\cpx{c}[1]&\g\h
\end{smallmatrix}\right)$, where $\cpx{c}$ is a homotopy for $\g\h\f\sim 0$.

\item[$(5)$] Assume that $\cpx{r}$ and $\cpx{s}$ are isomorphisms and $\g$ is a radical morphism in $\Kc$. If there exists a morphism $\cpx{\phi}:\cone(\f)\ra \cone(\g)$  of complexes such that both squares in $(\divideontimes)$ commute in $\Kc$, then $\cpx{\phi}$ is an isomorphism in $\Kc$.
\end{enumerate}
\end{lemma}
\begin{proof}


(1) Assume that the left square in $(\divideontimes)$ commutes. Then
$$\left(\begin{matrix}
\cpx{b}\\
\cpx{d}
\end{matrix}\right)= \left(\begin{matrix}
\cpx{a}& \cpx{b}\\
\cpx{c}& \cpx{d}
\end{matrix}\right)\circ \theta(\f)\sim \theta(\g)\circ \cpx{r}=\left(\begin{matrix}
0\\
\cpx{r}
\end{matrix}\right).$$
Let $\left(\begin{smallmatrix}
\cpx{m}\\
\cpx{n}
\end{smallmatrix}\right): \Y\ra \cone(\g)[-1]$ be a homotopy of $\left(\begin{smallmatrix}
\cpx{b}\\
\cpx{d}
\end{smallmatrix}\right)-\left(\begin{smallmatrix}
0\\
\cpx{r}
\end{smallmatrix}\right)$, that is,
\[\left(\begin{matrix}
-d_U^i&0\\
g^i&d_V^{i-1}
\end{matrix}\right)\left(\begin{matrix}
m^i\\
n^i
\end{matrix}\right)+\left(\begin{matrix}
m^{i+1}\\
n^{i+1}
\end{matrix}\right)d_Y^i=\left(\begin{matrix}
b^i\\
d^i-r^i
\end{matrix}\right)\]
 for all $i$. Then $-d_U^im^i+m^{i+1}d_Y^i=b^i$. This gives a graded morphism $\left(\begin{smallmatrix}
0 &-\cpx{m}\\
0&0
\end{smallmatrix}\right): \cone(\f)\ra \cone(\g)[-1]$ such that
$$\left(\begin{matrix}
0&-m^{i+1}\\
0&0
\end{matrix}\right)\left(\begin{matrix}
-d_X^{i+1}&0\\
f^{i+1}&d_Y^i
\end{matrix}\right)+\left(\begin{matrix}
-d_U^i&0\\
g^i&d_V^{i-1}
\end{matrix}\right)\left(\begin{matrix}
0&-m^i\\
0&0
\end{matrix}\right)+\left(\begin{matrix}
a^i&b^i\\
c^i&d^i
\end{matrix}\right)=\left(\begin{matrix}
a^i-m^{i+1}f^{i+1}&0\\
c^i&d^i-g^i m^i
\end{matrix}\right).$$
This means
$$\left(\begin{matrix}
\cpx{a}& \cpx{b}\\
\cpx{c}& \cpx{d}
\end{matrix}\right)\sim\left(\begin{matrix}
\cpx{a}-(\cpx{m}[1])(\f[1])&0\\
\cpx{c}&\cpx{d}-\g\cpx{m}
\end{matrix}\right): \cone(\f)\lra \cone(\g).$$
 The case in which the right square commutes can be shown similarly.

(2) This follows by definition.

(3)  Assume that there is a morphism $\left(\begin{smallmatrix}
\cpx{a}& \cpx{b}\\
\cpx{c}& \cpx{d}
\end{smallmatrix}\right): \cone(\f)\ra \cone(\g)$ of complexes such that the squares in $(\divideontimes)$ commute. By (1), we may assume $\cpx{b}=0$. The commutativity of the left square implies that
$$\left(\begin{matrix}
0\\
\cpx{d}
\end{matrix}\right)= \left(\begin{matrix}
\cpx{a}& 0\\
\cpx{c}& \cpx{d}
\end{matrix}\right)\circ \theta(\f)\sim \theta(\g)\circ \cpx{r}=\left(\begin{matrix}
0\\
\cpx{r}
\end{matrix}\right).$$
Thus $\theta(\g) (\cpx{d}-\cpx{r})=\left(\begin{smallmatrix}
0\\
\cpx{d}-\cpx{r}
\end{smallmatrix}\right)\sim 0$. Hence there is a morphism $\h: \Y\ra \U$ such that $\cpx{d}-\cpx{r}\sim \g\h$. Similarly, by the commutativity of the right square, there is a morphism $\cpx{m}: \Y[1]\ra \U[1]$ such that $\cpx{a}-\cpx{s}[1]\sim \cpx{m}(\f[1])$. Let $\cpx{w}: \X[1]\ra (\U[1])[-1]$ be a homotopy of $\cpx{a}-\cpx{s}[1]-\cpx{m}(\f[1])$, that is, $w^{i+1}(-d_X^{i+1})+(-d_U^i)w^i=a^i-s^{i+1}-m^if^{i+1}$ for all $i$. This gives rise to a graded morphism $\left(\begin{smallmatrix}
-\cpx{w} &-\cpx{m}[-1]\\
0&0
\end{smallmatrix}\right): \cone(\f)\ra \cone(\g)[-1]$ such that
$$\left(\begin{matrix}
-w^{i+1}&-m^i\\
0&0
\end{matrix}\right)\left(\begin{matrix}
-d_X^{i+1}&0\\
f^{i+1}&d_Y^i
\end{matrix}\right)+\left(\begin{matrix}
-d_U^i&0\\
g^i&d_V^{i-1}
\end{matrix}\right)\left(\begin{matrix}
-w^i&-m^{i-1}\\
0&0
\end{matrix}\right)+\left(\begin{matrix}
a^i&0\\
c^i&d^i
\end{matrix}\right)=\left(\begin{matrix}
s^{i+1}&0\\
-g^iw^i+c^i&-g^im^{i-1}+d^i
\end{matrix}\right).$$
Thus $\left(\begin{smallmatrix}
\cpx{a}& 0\\
\cpx{c}& \cpx{d}
\end{smallmatrix}\right)\sim\left(\begin{smallmatrix}
\cpx{s}[1]&0\\
-\g\cpx{w}+\cpx{c}&-\g(\cpx{m}[-1])+\cpx{d}
\end{smallmatrix}\right)$. So we may assume $\cpx{a}=\cpx{s}[1]$. By (2), $\cpx{d}\f\sim \g(\cpx{a}[-1])=\g\cpx{s}$. Thanks to $\cpx{d}-\cpx{r}\sim \g\h$, we have
$$\g\h\f\sim (\cpx{d}-\cpx{r})\f=\cpx{d}\f-\cpx{r}\f\sim \g\cpx{s}-\cpx{r}\f.$$

Conversely, assume that there is a morphism $\h: \Y\ra \U$ such that $\g\h\f\sim \g\cpx{s}-\cpx{r}\f.$ Let $\cpx{d}:=\cpx{r}+\g\h: \Y\ra \V$. Then $\g\cpx{s}\sim \cpx{r}\f+\g\h\f=(\cpx{r}+\g\h)\f=\cpx{d}\f$. Let $\cpx{c}$ be a homotopy of $\cpx{d}\f-\g\cpx{s}$. By (2),
$$\cpx{\phi}:=\left(\begin{matrix}
\cpx{s}[1]&0 \\
\cpx{c}[1]&\cpx{d}
\end{matrix}\right): \cone(\f)\lra \cone(\g)$$
 is a morphism of complexes. Clearly, $\pi(\g)\cpx{\phi}=(\cpx{s}[1]) \pi(\f)$. Since $\cpx{\phi} \theta(\f)=\left(\begin{smallmatrix}
0\\
\cpx{d}
\end{smallmatrix}\right)=\theta(\g) \cpx{d}$ and $\theta(\g) \g\sim 0$, we get
$$\cpx{\phi}\theta(\f)-\theta(\g) \cpx{r}=\theta(\g) \cpx{d}-\theta(\g) \cpx{r}=\theta(\g)(\cpx{d}-\cpx{r})= \theta(\g) \g\h\sim 0.$$

(4) Since both $\cpx{\phi}, \cpx{\psi}: \cone(\f)\ra \cone(\g)$ make the two squares in $(\divideontimes)$ commute, we have the following commutative diagram
\[\begin{array}{c}\xymatrix@R=1.5em{
\X\ar[r]^{\f}&\Y\ar[r]^-{\theta(\f)}\ar[d]^{0}& \cone(\f)\ar[r]^{\pi(\f)}
\ar[d]^{\cpx{\phi}-\cpx{\psi}}&\X[1]\ar[d]^{0} \\
\U\ar[r]^{\g}&\V\ar[r]^-{\theta(\g)}& \cone(\g)\ar[r]^{\pi(\g)}&\U[1].
}\end{array}\]
Now, the desired result follows from (3).

(5) By (3), there exists a morphism $\h: \Y\ra \U$ of complexes such that $\g\h\f\sim \g\cpx{s}-\cpx{r}\f$ with $\cpx{c}$ a homotopy, and
$$\cpx{\phi}=\left(\begin{matrix}
\cpx{s}[1]&0 \\
\cpx{c}[1]&\cpx{r}+\g\h
\end{matrix}\right): \cone(\f)\lra \cone(\g).$$
Then the following diagram
 \[\begin{array}{c}\xymatrix@R=1.5em{
\X\ar[r]^{\f}\ar[d]^{\s}&\Y\ar[r]^-{\theta(\f)}\ar[d]^{\cpx{r}+\g\h}& \cone(\f)\ar[r]^{\pi(\f)}
\ar[d]^{\cpx{\phi}}&\X[1]\ar[d]^{\cpx{s}[1]} \\
\U\ar[r]^{\g}&\V\ar[r]^-{\theta(\g)}& \cone(\g)\ar[r]^{\pi(\g)}&\U[1]
}\end{array}\]
is commutative in $\Kc$. Since $\cpx{r}$ and $\cpx{s}$ are isomorphisms and $\g$ is a radical morphism in $\Kc$, both $\s$ and $\cpx{r}+\g\h$ are isomorphisms in $\Kc$. It follows from the above diagram that $\cpx{\phi}$ is an isomorphism in $\Kc$.
\end{proof}

\subsection{Commutative Frobenius algebras}\label{CFA}

Throughout this subsection, $A$ denotes a finite-dimensional local commutative Frobenius algebra over a field $k$, and $\omega:A\ra k$ is a Frobenius form of $A$.

Let $U, V$ be finite-dimensional $k$-spaces. Assume that there is a non-degenerate $k$-bilinear form $\lan -,-\ran: U\times V\ra k$. Let $W$ be a subspace of $U$. Define
$$W^{\perp}=\{v\in V \mid \lan w, v\ran =0, \forall w\in W\}.$$
Then ${}^{\perp}(W^{\perp})=W$.
Assume that there is another non-degenerate $k$-bilinear form $\lan -,-\ran: V\times U\ra k$ such that $\lan u, v\ran=\alpha\lan v, u\ran$ for $u\in U, v\in V$, where $\alpha\in k \setminus \{0\}$ is a fixed element. Then $W^{\perp} = {}^{\perp}W$ with respect to the two bilinear forms, respectively. In this case, we just write $W^{\perp}$.

For an $A$-module  endomorphism $f$ of a free $A$-module $X$ with finite rank, we define $\tr(f)$ to be the trace of the matrix of $f$ with respect to a chosen basis of $X$. If $X=0$, then we define $\tr(f)=0$. For two  free $A$-modules $X$ and $Y$ with finite rank, since $A$ is commutative, one has $\tr(fg)=\tr(gf)$ for all $f\in\Hom_A(X,Y)$ and $g\in\Hom_A(Y,X)$. Clearly, the  function $\tr$ is $k$-linear and  independent of the choice of basis. Since the algebra $A$ is local, every finitely generated projective $A$-module is isomorphic to $A^m$ for some $m\geq 0$. Thus the function $\tr$ is defined for $\End_A(X)$ for every finitely generated projective $A$-module $X$.

\medskip
Let $\X$ and $\Y$ be two complexes in $\Kb{A}$. The Hom-complex $\Hom^{\bullet}(\X, \Y)$ of $\X$ and $\Y$ is defined as follows. For $n\in \mathbb{Z}$, its $n$th term is
$$\Hom^n(\X, \Y):=\prod_{p\in \mathbb{Z}}\Hom_A(X^p, Y^{p+n}),$$
and the differential is given by
$$d^n:  \Hom^n(\X, \Y)\lra \Hom^{n+1}(\X, \Y),$$
$$ (f^p: X^p\ra Y^{p+n})_{p\in \mathbb{Z}}\longmapsto \big(d_Y^{n+p}f^p+(-1)^{n+1}f^{p+1}d_X^p: X^p\ra Y^{p+n+1}\big)_{p\in \mathbb{Z}}.$$
As $\X$ and $\Y$ are bounded complexes, $\Hom^n(\X, \Y)$ is a finite-dimensional $k$-space for all $n\in \mathbb{Z}$. For each $n\in \mathbb{Z}$, define
$$Z^n(\Hom^{\bullet}(\X, \Y)):=\ker(d^n),\;\;\; B^n(\Hom^{\bullet}(\X, \Y)):=\im(d^{n-1}).$$


The following result may be known. For the convenience of the reader, we include here a proof.
\begin{lemma}\label{form}
Keep the notations above, and let $n\in\mathbb{Z}$. Then
 $$ \lan -,-\ran: \Hom^n(\X, \Y)\times \Hom^{-n}(\Y, \X)   \lra  k, $$
$$(\f,\g)  \longmapsto \sum_{p\in \mathbb{Z}}(-1)^p\omega(\tr(f^{p-n}g^p)) \; \mbox{ for } \f\in  \prod_{p\in \mathbb{Z}}\Hom_A(X^p, Y^{p+n}), \, \g\in \prod_{p\in \mathbb{Z}}\Hom_A(Y^p, X^{p-n})$$
is a non-degenerate $k$-bilinear form, satisfying the following properties.

$(1)$ $\lan \f, \g \ran=(-1)^n\lan \g, \f \ran$ for  $\f\in \Hom^n(\X, \Y)$ and $\g\in \Hom^{-n}(\Y, \X)$.

$(2)$ $\lan d^{n-1}\cpx{h}, \cpx{z}\ran =(-1)^{n}\lan \cpx{h}, d^{-n}\cpx{z}\ran$ for $\cpx{h}\in \Hom^{n-1}(\X, \Y)$ and $\cpx{z}\in \Hom^{-n}(\Y, \X)$.

$(3)$ Let $(-)^\perp$ be the orthogonal subspace with respect to the bilinear form $\lan -,-\ran$. Then
 \[B^n(\Hom^{\bullet}(\X, \Y))^{\perp}=Z^{-n}(\Hom^{\bullet}(\Y, \X)) \mbox{ and }
 Z^{-n}(\Hom^{\bullet}(\Y, \X))^{\perp}=B^n(\Hom^{\bullet}(\X, \Y)).\]
\end{lemma}
\begin{proof}
Since $\X$ and $\Y$ are bounded,  the summation in the definition of $\lan \f, \g \ran$ is a finite sum. It is evidently $k$-bilinear. Now if $\g\in\Hom^{-n}(\Y,\X)$ is nonzero, then there is some $p\in\Z$ such that $g^p: Y^p\ra X^{p-n}$ is nonzero. There are positive integers $r,s$ such that $Y^p\simeq A^r$ and $X^{p-n}\simeq A^s$, and $g^p$ can be viewed as a $s\times r$-matrix over $A$. Suppose the $(i,j)$-entry $a$ of $g^p$ is nonzero. Since $\omega$ is a Frobenius form of $A$, there is an element $b\in A$ such that $\omega(ba)\neq 0$. Let $h$ be a $r\times s$-matrix over $A$ with $(j,i)$-entry $b$ and other entries zero. Then $\omega\tr(hg^p)=\omega(ba)\neq 0$. The matrix $h$ can be viewed as an $A$-module morphism from $X^{p-n}$ to $Y^p$. Let $\f\in\Hom^n(\X,\Y)$ be such that $f^j=0$ for all $j\neq p-n$ and  $f^{p-n}=h$. Then we obtain that $\lan \f, \g\ran=(-1)^p\omega(\tr(hg^p))\neq 0$. Similarly, for  $\f\ne 0$, we can find $\g$ such that $\lan\f,\g\ran\ne 0$. Therefore, the bilinear form $\lan-,-\ran$ is non-degenerate.

To see (1), let $\f\in\Hom^n(\X,\Y)$ and $\g\in\Hom^{-n}(\Y,\X)$. Setting $q=p-n$, we obtain
\[
\begin{aligned}
\lan\f,\g\ran
&=\sum_{p\in\mathbb Z}(-1)^p\omega\bigl(\tr(f^{p-n}g^p)\bigr) =\sum_{p\in\mathbb Z}(-1)^p\omega\bigl(\tr(g^pf^{p-n})\bigr)\\
&=\sum_{q\in\mathbb Z}(-1)^{q+n}\omega\bigl(\tr(g^{q+n}f^q)\bigr)
=(-1)^n\lan\g,\f\ran.
\end{aligned}
\]

We prove (2). The relevant components of the differentials are
\[
(d^{n-1}\cpx{h})^{p-n}=d_Y^{p-1}h^{p-n}+(-1)^nh^{p-n+1}d_X^{p-n} \text{ and }
(d^{-n}\cpx{z})^p=d_X^{p-n}z^p+(-1)^{1-n}z^{p+1}d_Y^p.
\]
Consequently,
\[
\begin{aligned}
\lan d^{n-1}\cpx{h},\cpx{z}\ran
&=\sum_{p\in\mathbb Z}(-1)^p\omega\Bigl(\tr\bigl(
  (d_Y^{p-1}h^{p-n}+(-1)^nh^{p-n+1}d_X^{p-n})z^p\bigr)\Bigr)\\
&=\sum_{p\in\mathbb Z}(-1)^{p+1}\omega\bigl(\tr(h^{p-n+1}z^{p+1}d_Y^p)\bigr)
  +\sum_{p\in\mathbb Z}(-1)^{p+n}\omega\bigl(\tr(h^{p-n+1}d_X^{p-n}z^p)\bigr)\\
&=(-1)^n\sum_{p\in\mathbb Z}(-1)^p\omega\Bigl(\tr\bigl(
  h^{p-n+1}(d_X^{p-n}z^p+(-1)^{1-n}z^{p+1}d_Y^p)\bigr)\Bigr)\\
&=(-1)^n\lan\cpx{h},d^{-n}\cpx{z}\ran.
\end{aligned}
\]
Here, we have shifted the index by one in the first summand of the second line, and then used $\tr(MN)=\tr(NM)$. This proves (2).

Finally, we show (3). Let $\cpx{z}\in\Hom^{-n}(\Y,\X)$. By (2), the element $\cpx{z}$ is orthogonal to $B^n(\Hom^{\bullet}(\X,\Y))$ if and only if $\lan\cpx{h},d^{-n}\cpx{z}\ran=0$ for every $\cpx{h}\in\Hom^{n-1}(\X,\Y)$. This is equivalent to saying that $d^{-n}\cpx{z}=0$ because $\lan-,-\ran$ is non-degenerate.  Hence $B^n(\Hom^{\bullet}(\X,\Y))^{\perp}=Z^{-n}(\Hom^{\bullet}(\Y,\X))$. Further, taking orthogonal complements, we get
\[B^n(\Hom^{\bullet}(\X,\Y))=(B^n(\Hom^{\bullet}(\X,\Y))^{\perp})^{\perp} = Z^{-n}(\Hom^{\bullet}(\Y,\X))^{\perp}.\qedhere\]
\end{proof}

\section{A characterization of standard derived equivalences}\label{aut}
In this section, we study derived equivalences between finite-dimensional algebras and prove Theorem \ref{t0}.

Let $A$ and $B$ be finite-dimensional algebras over a common field, and let $F: \Db{A}\ra \Db{B}$ be a derived equivalence. It was proved in \cite{R} that  there exists a standard derived equivalence $\w{F}: \Db{A}\ra \Db{B}$ such that $F$ and $\w{F}$ agree on $A\proj$ and $F(\X)\simeq \w{F}(\X)$ for every object $\X$ of $\Db{A}$. Determining whether $F$ is standard is equivalent to understanding the natural-isomorphism class of the difference $\w{F}^{-1}F: \Db{A}\ra \Db{A}$. Our aim is to find a particularly convenient representative of this class.

We proceed in two steps. More precisely, a derived equivalence $F: \Db{A}\ra \Db{B}$ is a pair $(F, \xi): \Db{A}\ra \Db{B}$ consisting of an equivalence $F$ and a connecting isomorphism $\xi: F\circ [1]\ra [1]\circ F$ (the definition of a triangle functor is recalled below). First, we treat the connecting isomorphism. We show that, up to natural isomorphism of triangle functors, the connecting isomorphism can be taken to be the identity (see Proposition \ref{id}). Second, we consider a special class of triangles---the truncated triangles---and show that, up to natural isomorphism of triangle functors, the difference $\w{F}^{-1}F$ fixes all objects, all subcategories $A\md[i]$ for $i\in \Z$, and all truncated triangles. Both steps together will yield Theorem \ref{t0}.

Let $\C$ and $\A$ be triangulated categories with suspension functors $\Sigma$ and $\Sigma'$, respectively. A \emph{triangle functor} from $\C$ to $\A$ is a pair $(F,\phi)$ consisting of an additive functor $F:\C\ra\A$ and a natural isomorphism $\phi:F\Sigma\ra\Sigma'F$. These data are required to preserve triangles: whenever
$X\raf{f}Y\raf{g}Z\raf{h}\Sigma(X)$
is a triangle in $\C$, the sequence
$F(X)\lraf{F(f)}F(Y)\lraf{F(g)}F(Z)\lraf{\phi_XF(h)}\Sigma' (F(X))$
is a triangle in $\A$.   The natural isomorphism $\phi$ is called a \emph{connecting isomorphism} of $F$. If $F$ is an equivalence, then $(F,\phi)$ is called a \emph{triangle equivalence}. In this case, the categories $\C$ and $\A$ are said to be \emph{triangle equivalent}. Given another triangle functor $(G,\psi):\C\ra\A$, a natural transformation $\mu:(F,\phi)\ra(G,\psi)$ of triangle functors is a natural transformation $\mu:F\ra G$ compatible with the connecting isomorphisms; explicitly, the following diagram must commute:
  \[\begin{array}{c}\xymatrix@R=1.5em{
F\Sigma\ar[r]^{\phi}\ar[d]_{\mu\Sigma} &\Sigma' F\ar[d]^{\Sigma'\mu}\\
G\Sigma\ar[r]^{\psi}&\Sigma' G.
}\end{array}\]


We next recall a useful modification of a triangle functor. Let $(F,\phi):\C\ra\A$ be a triangle functor, let $\B$ be a full subcategory of $\C$, and choose a family of isomorphisms
$$\mu_{\B}=\{\mu_{X}: F(X)\ra F(X)\;|\; X\in \B\}$$
in $\A$. Extend this family to all objects of $\C$ by setting $\mu_Z:=\id_{F(Z)}$ whenever $Z\notin\B$. Define a functor $G:\C\ra\A$ by
$$G(X)=F(X),\quad G(f)=\mu_YF(f)\mu_X^{-1}$$
for every object $X$ and morphism $f:X\ra Y$ in $\C$. Its connecting isomorphism $\psi:G\Sigma\ra\Sigma'G$ is given by
$$\psi_X=\Sigma'(\mu_X)\phi_X\mu_{\Sigma(X)}^{-1}:G(\Sigma(X))\lra\Sigma'(G(X)).$$
Then $(G,\psi)$ is a triangle functor, and the family $\mu$ defines a natural isomorphism $(F,\phi)\ra(G,\psi)$ of triangle functors. We call $(G,\psi)$ the {\it conjugate functor} of $(F,\phi)$ by $\mu_{\B}$.

A category $\C$ is {\it essentially small} if its isomorphism classes of objects form a set. For an object $X\in\C$, we denote its isomorphism class by $[X]$.

\begin{lemma}\label{connect}
Let $(\C, \Sigma)$ and $(\A, \Sigma')$ be two essentially small triangulated categories. Let $(F, \phi): \C\ra \A$ be a triangle functor such that $F(\Sigma(X))=\Sigma' (F(X))$ for every object $X$ of $\C$. Then there is a triangle functor $(F', \id): \C\ra \A$, with connecting isomorphism $\id: F'\Sigma \ra \Sigma' F'$, such that $(F, \phi)\simeq (F', \id)$ as triangle functors.
\end{lemma}
\begin{proof}
Since the group generated by $\Sigma$ acts on the set $[\C]$ of all isomorphism classes of objects of $\C$, we can choose a subset $\mathcal{U}$ of $[\C]$ consisting of representatives of the orbits of this action; that is, for every object $X$ of $\C$, there exists an $n\in \mathbb{Z}$ such that $[X]\in \Sigma^n\mathcal{U}$.

Let $\mu_X=\id_{F(X)}$ for all $X$ with $[X]\in \mathcal{U}$. We will construct isomorphisms $\mu_X: F(X)\ra F(X)$ with $[X]\in \Sigma^n\mathcal{U}$ by induction on $n\leq 0$. Assume that we constructed $\mu_X$ for all objects $X$ with $[X]\in \bigcup_{n\leq i\leq 0}\Sigma^i(\mathcal{U})$. Let $X$ be an object of $\C$ with $[X]\in \Sigma^{n-1}\mathcal{U}$. Then $[\Sigma(X)]\in \Sigma^n\mathcal{U}$. Let $\mu_X:=\Sigma'^{-1}(\mu_{\Sigma(X)}\circ \phi_{X}^{-1})$. We have the following commutative diagram
  \[\begin{array}{c}\xymatrix{
F(\Sigma(X))\ar[r]^{\phi_{X}}\ar[d]_{\mu_{\Sigma(X)}} &\Sigma' (F(X))\ar[d]^{\Sigma'(\mu_X)}\\
F'(\Sigma(X))\ar[r]^{\id}&\Sigma' (F'(X)).
}\end{array}\]
 Thus we can construct $\mu$ on all objects $X$ with $[X]\in \bigcup_{-\infty<i\leq 0}\Sigma^i\mathcal{U}$. The construction for objects $X$ with $[X]\in \bigcup_{0\leq i<+\infty}\Sigma^i\mathcal{U}$ is similar. Denote the resulting family of automorphisms by $\mu_{\C}$. Let $(F', \psi): \C\ra \A$ be the conjugate functor of $(F, \phi)$ by $\mu_{\C}$. By construction, $\psi=\id$.
\end{proof}


Let $A$ be an Artin algebra. Following \cite{CY}, a triangle functor $(G, \phi): \Db{A}\ra \Db{A}$ (respectively, $(F, \xi): \Kb{A}\ra \Kb{A}$) is called a {\it pseudo-identity} if the following two conditions hold:
\begin{enumerate}
\item $G$ fixes $A\md[i]$ (respectively, $F$ fixes $A\proj[i]$) for each $i\in \Z$;
\item $G$ fixes all objects of $\Db{A}$ (respectively, $F$ fixes all objects of $\Kb{A}$).
\end{enumerate}
Here, we say that a triangle functor $H: \A\ra \A$ \emph{fixes} a full subcategory $\B$ (respectively, a triangle $\Delta$ or an ideal $\mathcal{I}$) of $\A$ if $H$ fixes all objects and morphisms in $\B$ (respectively, all objects and morphisms in $\Delta$ or all morphisms in $\mathcal{I}$).

\medskip
The following results are due to \cite[Proposition 5.8]{CY}, motivated by \cite[Corollary 3.5]{R}. Note that the proof of (2) is similar to that of (1).

\begin{lemma}\label{group}
Suppose that $A$ and $B$ are finite-dimensional algebras over a common field.
\begin{enumerate}
\item[$(1)$] Let $(F, \xi): \Db{A}\ra \Db{B}$ be a derived equivalence. Then there is a standard derived equivalence $(G, \psi): \Db{A}\ra \Db{B}$ and a pseudo-identity $(I, \phi): \Db{A}\ra \Db{A}$ such that $(F, \xi)\simeq(G, \psi)\circ(I, \phi)$ as triangle functors.
\item[$(2)$] Let $(F, \xi): \Kb{A}\ra \Kb{B}$ be a derived equivalence. Then there is a standard derived equivalence $(G, \psi): \Kb{A}\ra \Kb{B}$ and a pseudo-identity $(I, \phi): \Kb{A}\ra \Kb{A}$ such that $(F, \xi)\simeq (G, \psi)\circ (I, \phi)$ as triangle functors.
\end{enumerate}
\end{lemma}

\begin{prop}\label{id}
Suppose that $A$ and $B$ are finite-dimensional algebras over a common field.
\begin{enumerate}
\item[$(1)$] Every derived equivalence $(F, \xi): \Db{A}\ra \Db{B}$ is naturally isomorphic to a derived equivalence $(F', \id): \Db{A}\ra \Db{B}$. Moreover, if $(F, \xi)$ is a pseudo-identity, then so is $(F', \id)$.
\item[$(2)$] Every derived equivalence $(F, \xi): \Kb{A}\ra \Kb{B}$ is naturally isomorphic to a derived equivalence $(F', \id): \Kb{A}\ra \Kb{B}$. Moreover, if $(F, \xi)$ is a pseudo-identity, then so is $(F', \id)$.
\end{enumerate}
\end{prop}
\begin{proof}
  (1) Every pseudo-identity $(H, \phi): \Db{A}\ra \Db{A}$ satisfies $H(\X[1])=H(\X)[1]$ for every object $\X$ of $\Db{A}$. Hence, by Lemma \ref{connect}, there is a natural isomorphism $\mu: (H, \phi)\simeq (H', \id)$ of triangle functors for some functor $H'$. By the construction in Lemma \ref{connect}, $H'$ fixes all objects. Note that, for each $i\in \Z$ and every object $\X$ of $A\md[i]$, the map $\phi_{\X}: \X[1]\ra \X[1]$ is given by multiplication by an invertible central element of $A$. The construction in Lemma \ref{connect} shows that $\mu_{\X}$ is also given by multiplication by an invertible central element of $A$; hence $H'$ fixes $A\md[i]$. Thus $(H', \id)$ is a pseudo-identity.

By definition, every standard equivalence $(G, \psi): \Db{A}\ra \Db{B}$ is naturally isomorphic to a derived tensor functor. Since every derived tensor functor commutes with suspension functors, $(G, \psi)\simeq (G', \id)$ as triangle functors for some functor $G'$. The assertion now follows from Lemma \ref{group}(1).

(2) The proof is similar to that of (1).
\end{proof}

\begin{defn}\label{def-truncated}
Let $A$ be an Artin algebra, and let $\X$ be a complex in $\Kb{A}$. The {\it truncated triangle of $\X$ in degree $i$} is defined to be the standard triangle
$$\Delta_i(\X):\;\; X^{\leq i}[-1]\lraf{\uu{d_X^i}} X^{\geq i+1}\lraf{\theta(\uu{d_X^i})} \X\lraf{\pi(\uu{d_X^i})} X^{\leq i}$$
in $\Kb{A}$, where $\uu{d_X^i}$ is the morphism with the $(i+1)$-component being $d_X^i$. We write $\Delta_L(\X)$ for $\Delta_m(\X)$, and $\Delta_R(\X)$ for $\Delta_{n-1}(\X)$, where $m$ is the smallest integer such that $X^m\neq 0$ and $n$ is the largest integer such that $X^n\neq 0$. In this case, $\Delta_L(\X)$ and $\Delta_R(\X)$ are called the \emph{left and right} truncated triangles of $\X$, respectively.
For convenience, we write $\theta_L(\X):=\theta(\uu{d_X^m}), \pi_L(\X):=\pi(\uu{d_X^m}), \theta_R(\X):=\theta(\uu{d_X^{n-1}})$ and $\pi_R(\X):=\pi(\uu{d_X^{n-1}})$. In this case, the {\em length} $\ell(\X)$ of $\X$ is defined to be $n-m$.
\end{defn}

For our later considerations, we introduce the following ideals and subcategories of $\Kb{A}$ for an Artin algebra $A$. Let $i$ be an integer and $n$ a non-negative integer.
\begin{enumerate}
%  \item $\Stmap_i(A)$ is the ideal of $\Kb{A}$ consisting of the homotopy classes of morphisms $\f$ such that $f^t=0$ for all $t\neq i$;
  \item $\Zmap_i(A)$ is the ideal consisting of morphisms $\f:\X\ra\Y$ such that $f^t=0$ for all $t\neq i$ and $f^i$ factors through $d_Y^{i-1}$;
  \item $\K^{\ell\leq n}(A\proj)$ is the full subcategory of $\Kb{A}$ consisting of complexes of length at most $n$; and  \item $\K^{0,n}(A\proj)$ is the full subcategory of $\Kb{A}$ consisting of indecomposable complexes of length $n$ whose nonzero terms are concentrated in degrees $0,1,\ldots,n$.
\end{enumerate}

Let $\B$ be a full subcategory of $\Kb{A}$ closed under the shift functor $[1]$. A family of auto-isomorphisms $\lambda_{\B}=\{\lambda_{\X}:\X\ra \X\;|\; \X\in \B\}$ in $\Kb{A}$ is {\it $[1]$-compatible} if $\lambda_{\X[1]}=\lambda_{\X}[1]$ for all objects $\X$ in $\B$. A crucial property we shall use frequently below is that the conjugate functor $(F', \phi)$ of a pseudo-identity
$(F, \id): \Kb{A}\ra \Kb{A}$ by a $[1]$-compatible family $\lambda_{\B}$ also satisfies $\phi=\id: F'[1]\ra [1]F'$.

\begin{prop}\label{truncated}
Let $A$ be a finite-dimensional algebra, and let $(F, \phi): \Kb{A}\ra \Kb{A}$ be a pseudo-identity. Then there is a pseudo-identity $(G, \id): \Kb{A}\ra \Kb{A}$ such that $(G, \id)\simeq (F, \phi)$ as triangle functors and $G$ fixes all truncated triangles in $\Kb{A}$.
\end{prop}
\begin{proof}

By Proposition \ref{id}, we may assume $\phi=\id$.
We will construct, by induction on the length of complexes, a $[1]$-compatible family of isomorphisms  whose corresponding conjugate functor $(G,\id)$ fixes both left and right truncated triangles.

For complexes of length zero, take the identity isomorphisms. Since such complexes are direct sums of shifts of stalk complexes, the resulting functor $(F_0=F,\id)$ clearly fixes the left and right truncated triangles of those length zero complexes.

Suppose that $n > 0$ and that a $[1]$-compatible family $\lambda_{\K^{\ell\leq n-1}(A\proj)}$ has been constructed on objects of length at most $n-1$ such that the conjugate functor $(F_{n-1},\id)$ of $(F_0,\id)$ by this family fixes both left and right truncated triangles of complexes with length at most $n-1$. We extend the family to complexes of length $n$. Let $\K^{\ell=n}(A\proj)$ be the full subcategory of $\Kb{A}$ consisting of all complexes of length $n$. We first start with a complex $\cpx{X}$ in $\K^{0,n}(A\proj)$ and consider the left truncated triangle of $\cpx{X}$. Both the composition
\[X^0[-1]\lraf{\uu{d_X^0}}X^{\geq 1}\lraf{\pi_L(X^{\geq 1})}X^1[-1]\]
and $\pi_L(X^{\geq 1})$ are fixed by $F_{n-1}$ by induction, since the complexes have length less than $n$. It follows that $F_{n-1}$ also fixes $\uu{d_X^0}$. Thus there is an isomorphism $\lambda_{\X}^L: \X\ra \X$ in $\Kb{A}$ such that the following diagram
  \[\xymatrix{
F_{n-1}(\Delta_L(\X)): &X^0[-1]\ar[r]^{\uu{d_X^0}}\ar[d]^{1}&X^{\geq 1}\ar[d]^{1}\ar[rr]^-{F_{n-1}(\theta_L(\X))} &&\X\ar@{..>}[d]^{\lambda_{\X}^L}\ar[rr]^{F_{n-1}(\pi_L(\X))}&&X^0\ar[d]^{1} \\
\Delta_L(\X): &X^0[-1]\ar[r]^{\uu{d_X^0}}&X^{\geq 1}\ar[rr]^-{\theta_L(\X)} &&\X\ar[rr]^{\pi_L(\X)}&&X^0}\]
commutes in $\Kb{A}$. For any complex $\X\in\Kb{A}$ with length $n$, there is a unique integer $i$ such that $\X[i]\in\K^{0,n}(A\proj)$; define $\lambda_{\X}^L=\lambda_{\X[i]}^L[-i]$. Thus we get a $[1]$-compatible family $\lambda^L_{\K^{\ell=n}(A\proj)}$ of automorphisms on $\K^{\ell=n}(A\proj)$.  The conjugate functor $(F^L_n, \id)$ of $(F_{n-1}, \id)$ by $\lambda^L_{\K^{\ell=n}(A\proj)}$ is a pseudo-identity. Moreover, the right two commutative squares of the above diagram imply that $F^L_n$ fixes $\Delta_L(X)$ for all complexes $\cpx{X}$ with length $n$.


We now correct the right truncated triangle at the same induction stage. Let $\X$ be an object in $\K^{0,n}(A\proj)$. First, we claim that
$$(\star)\;\;\;F^L_n(\uu{d_X^{n-1}})=\uu{d_X^{n-1}},\;\;F^L_n(\theta_R(\X))=\theta_R(\X)\;\;\text{and}\;\;  F^L_n(\pi_R(\X))=\pi_R(\X)+\cpx{\phi}$$
for some $\cpx{\phi}\in \Zmap_1(A)(\X, X^{\leq n-1})$. Observe that
$$\uu{d_X^{n-1}}\theta_R(X^{\leq n-1}[-1])=d_X^{n-1}[-n]: X^{n-1}[-n]\lra X^n[-n].$$
Since $F^L_n$ fixes both $\uu{d_X^{n-1}}\theta_R(X^{\leq n-1}[-1])$ and $\theta_R(X^{\leq n-1}[-1])$, it fixes $\uu{d_X^{n-1}}$. This gives the first equality of $(\star)$. Since
$$\theta_R(\X)=\theta_L(X^{\geq 0})\theta_L(X^{\geq 1})\cdots \theta_L(X^{\geq n-2})\theta_L(X^{\geq n-1})$$
and $F^L_n$ fixes $\theta_L(X^{\geq i})$ for all $0\leq i\leq n-1$, the second equality of $(\star)$ follows.  To prove the last equality, consider the morphism of left truncated triangles
  \[\xymatrix{
\Delta_L(\X): &X^0[-1]\ar[rr]^-{\uu{d_X^0}}\ar[d]^1&&X^{\geq 1}\ar[d]^{\pi_R(X^{\geq 1})}\ar[rr]^-{\theta_L(\X)} &&\X\ar[d]^{\pi_R(\X)}\ar[rr]^{\pi_L(\X)}&&X^0\ar[d]^1 \\
\Delta_L(X^{\leq n-1}): &X^0[-1]\ar[rr]^-{\uu{d_{X^{\leq n-1}}^0}}&&X^{[1, n-1]}\ar[rr]^-{\theta_L(X^{\leq n-1})} &&X^{\leq n-1}\ar[rr]^{\pi_L(X^{\leq n-1})}&&X^0.}\]
Since $F^L_n$ fixes $\Delta_L(\X)$, $\Delta_L(X^{\leq n-1})$ and $\pi_R(X^{\geq 1})$, applying $F^L_n$ to the above diagram and using Lemma \ref{st}(4), we obtain the last equality of $(\star)$ for some $\cpx{\phi}\in \Zmap_1(A)(\X, X^{\leq n-1})$. By the definition of $\Zmap_1(A)$, we can represent $\cpx{\phi}$ so that $\phi^i=0$ unless $i=1$ and there is a module homomorphism $h: X^1\ra X^0$ such that $\phi^1=d_{X^{\leq n-1}}^0h$ and  $d_{X^{\leq n-1}}^0hd_X^0=0$. Note that  $d_{X^{\leq n-1}}^0=0$ if $n=1$.

For $n=1$, set $\cpx{\psi}:\X\lra \X$ to be the zero morphism, and for $n>1$, let
 $\cpx{\psi}:\cpx{X}\lra \cpx{X}$ be given by $\psi^1=d_X^0h$ and $\psi^i=0$ for all $i\neq 1$. Then $(\cpx{\psi})^2=0$, and $\mu_{\cpx{X}}=\id_{\X}+\cpx{\psi}$ is an isomorphism with $\mu_{\cpx{X}}^{-1}=\id_{\X}-\cpx{\psi}$.


In $\K^b(A\proj)$, by construction, we have $\cpx{\phi}=\pi_R(\cpx{X})\cpx{\psi}$ and $\cpx{\psi}\theta_R(\cpx{X})=0$.  For each $\X$ of length $n$, let $i$ be the unique integer such that $\X[i]\in\K^{0,n}(A\proj)$, and define $\mu_{\X}=\mu_{\X[i]}[-i]$. We get a $[1]$-compatible family of automorphisms in $\K^{\ell=n}(A\proj)$.
We denote by $(F_n,\id)$ the resulting functor obtained by conjugating $(F_n^L,\id)$ by this $[1]$-compatible family.

By the definition of conjugation, together with  the equalities in $(\star)$, we get
\[\begin{aligned}
  F_n(\pi_R(\X)) & =\id_{X^{\leq n-1}}F_{n}^L(\pi_R(\X))\mu_{\X}^{-1}\\
    & =(\pi_R(\X)+\cpx{\phi})(1-\cpx{\psi})\\
    & =(\pi_R(\X)+\pi_R(\X)\cpx{\psi})(1-\cpx{\psi})=\pi_R(\X).
\end{aligned}\]
Similarly, we have
\[F_n(\theta_R(\X))  =\mu_{\X}\, F_{n}^L(\theta_R(\X))\,\id_{X^{\geq n}}^{-1}
   =\mu_{\X}\,\theta_R(\X)=(\id_{\X}+\cpx{\psi})\theta_R(\X)=\theta_R(\X),\]
where the second equality follows from the second equality of $(\star)$. Since both $X^{\leq n-1}[-1]$ and $X^{\geq n}$ have length less than $n$, we also have $F_n(\uu{d_X^{n-1}})=F_{n}^L(\uu{d_X^{n-1}})=\uu{d_X^{n-1}}$. Hence  $F_n$ fixes $\Delta_R(\X)$ for all $\X\in\K^{0,n}(A\proj)$.

It remains to check that this last correction does not destroy the left truncated triangle fixed earlier by $F_n^L$. It is easy to check that
$\pi_L(\X)\cpx{\psi}=0$ and $\cpx{\psi}\theta_L(\X)=0$ in $\Kb{A}$; the second equality is represented by the null-homotopy $h: X^1\ra X^0$. Therefore
\[\mu_{\X}\theta_L(\X)=(\id_{\X}+\cpx{\psi})\theta_L(\X)=\theta_L(\X)\text{ and }\pi_L(\X)\mu_{\X}=\pi_L(\X)(\id_{\X}+\cpx{\psi})=\pi_L(\X).\]
This shows that $F_n$ still fixes $\Delta_L(\X)$ for $\X\in\K^{0,n}(A\proj)$.  Since the connecting isomorphism $[1]F_n\ra F_n[1]$ is still $\id$, we conclude that $F_n$ fixes $\Delta_L(\X)$ and $\Delta_R(\X)$ for all complexes $\X$ with length $n$.

Altogether, we have constructed a $[1]$-compatible family of isomorphisms $\mu_{\X}\lambda_{\X}^L$ on $\K^{\ell=n}(A\proj)$ such that $(F_n,\id)$ is the conjugate functor of $(F_{n-1},\id)$ by this family. Denote this family by $\lambda_{\K^{\ell=n}(A\proj)}$, and let $\lambda_{\K^{\ell\leq n}(A\proj)}$ be the union of $\lambda_{\K^{\ell\leq n-1}(A\proj)}$ and $\lambda_{\K^{\ell=n}(A\proj)}$. Then $(F_n,\id)$ is the conjugate functor of $(F_0,\id)$ by $\lambda_{\K^{\ell\leq n}(A\proj)}$.

Passing to the union of the $[1]$-compatible families over all lengths gives a pseudo-identity $(G,\id)$, naturally isomorphic to $(F, \id)$, which fixes  $\Delta_L(\X)$ and $\Delta_R(\X)$ for all complexes $\X$.

We now show that $(G, \id)$ fixes the truncated triangles $\Delta_t(\X)$ for all $\X$ and all $t$.  Assume that $X^i=0$ for all $i<m$ and all $i>n$, and that $X^m\neq 0\neq X^n$. We may assume $m\leq t<n$.  Consider $\uu{d_X^t}: X^{\leq t}[-1]\ra X^{\geq t+1}$. One can directly check that
$$\pi_L(X^{\geq t+1})\circ \uu{d_X^t}\circ \theta_R(X^{\leq t}[-1])=d_X^t[-t-1]: X^t[-t-1]\lra X^{t+1}[-t-1].$$
 Since $G$ fixes $d_X^t[-t-1]$, $\theta_R(X^{\leq t}[-1])$ and $\pi_L(X^{\geq t+1})$, it fixes $\uu{d_X^t}$. Moreover,
$$\theta(\uu{d_X^t})=\theta_L(X^{\geq m})\cdots \theta_L(X^{\geq t-1})\theta_L(X^{\geq t}),\;\;\; \pi(\uu{d_X^t})=\pi_R(X^{\leq t+1})\cdots \pi_R(X^{\leq n-1})\pi_R(X^{\leq n}).$$
Since $G$ fixes $\theta_L(\Y)$ and $\pi_R(\Y)$ for all objects $\Y$ in $\Kb{A}$, $G$ fixes $\theta(\uu{d_X^t})$ and $\pi(\uu{d_X^t})$. So $G$ fixes $\Delta_t(\X)$.
\end{proof}

Now we are ready to prove Theorem \ref{t0}. Note that a pseudo-identity $\Db{A}\ra \Db{A}$ or $\Kb{A}\ra \Kb{A}$ is standard if and only if it is naturally isomorphic to the identity triangle functor.

\begin{proof} [{\bf Proof of Theorem \ref{t0}}]
It is known that every derived auto-equivalence of $\Db{A}$ restricts to a derived auto-equivalence of $\Kb{A}$. Let $\Out(\Db{A})$ (respectively, $\Out(\Kb{A})$) be the group of natural-isomorphism classes of derived auto-equivalences of $\Db{A}$ (respectively, $\Kb{A}$). Let $\res: \Out(\Db{A})\ra \Out(\Kb{A})$ be the restriction map. Then $\res$ is a group isomorphism \cite[Proposition, p. 658]{C1}. Hence a derived auto-equivalence $F: \Db{A}\ra \Db{A}$ is standard if and only if $\res(F): \Kb{A}\ra \Kb{A}$ is standard. Therefore, Rickard's question holds for $A$ if and only if every derived auto-equivalence of $\Kb{A}$ is standard. By Lemma \ref{group}, this is equivalent to every pseudo-identity of $\Kb{A}$ being standard. By Proposition \ref{truncated}, the latter is equivalent to the statement: every pseudo-identity $(F, \id): \Kb{A}\ra \Kb{A}$ that fixes all truncated triangles is naturally isomorphic to the identity triangle functor.
\end{proof}

To make Theorem \ref{t0} useful, we need to understand pseudo-identities that fix all truncated triangles. First, we establish some properties of such pseudo-identities.

Let $(F, \id): \Kb{A}\ra \Kb{A}$ be a pseudo-identity that fixes all truncated triangles. For a morphism $\f: \X\ra \Y$ in $\Kb{A}$ and an integer $i\in \Z$, consider the following morphism of triangles:
  \[(*)\;\;\;\xymatrix@C=15mm{
\Delta_i(\X): &X^{\leq i}[-1]\ar[r]^-{\uu{d_X^i}}\ar[d]^{f^{\leq i}[-1]}&X^{\geq i+1}\ar[d]^{f^{\geq i+1}}\ar[r]^-{\theta(\uu{d_X^i})} &\X\ar[d]^{\f}\ar[r]^{\pi(\uu{d_X^i})}&X^{\leq i}\ar[d]^{f^{\leq i}} \\
\Delta_i(\Y): &Y^{\leq i}[-1]\ar[r]^-{\uu{d_Y^i}}&Y^{\geq i+1}\ar[r]^-{\theta(\uu{d_Y^i})} &\Y\ar[r]^{\pi(\uu{d_Y^i})}&Y^{\leq i}.}\]
Define a graded morphism
$$\chi_{i,F}(\f): \X\lra \Y$$
by $\chi_{i,F}(\f)^{\leq i}=F(f^{\leq i})$ and $\chi_{i,F}(\f)^{\geq i+1}=F(f^{\geq i+1})$. Applying $F$ to the left square of $(*)$ shows that $\chi_{i,F}(\f)$ is a morphism of complexes. Moreover, $\chi_{i,F}(\f)$ makes the following diagram
   \[(**)\;\;\xymatrix@C=15mm{
\Delta_i(\X): &X^{\leq i}[-1]\ar[r]^-{\uu{d_X^i}}\ar[d]^{F(f^{\leq i})[-1]}&X^{\geq i+1}\ar[d]^{F(f^{\geq i+1})}\ar[r]^-{\theta(\uu{d_X^i})} &\X\ar[d]^{\chi_{i,F}(\f)}\ar[r]^{\pi(\uu{d_X^i})}&X^{\leq i}\ar[d]^{F(f^{\leq i})} \\
\Delta_i(\Y): &Y^{\leq i}[-1]\ar[r]^-{\uu{d_Y^i}}&Y^{\geq i+1}\ar[r]^-{\theta(\uu{d_Y^i})} &\Y\ar[r]^{\pi(\uu{d_Y^i})}&Y^{\leq i}.}\]
commutes in $\Kb{A}$.

\begin{lemma}[Reduction Lemma]\label{coro}
Let $A$ be a finite-dimensional algebra, and let $(F, \id): \Kb{A}\ra \Kb{A}$ be a pseudo-identity that fixes all truncated triangles.
\begin{enumerate}
\item[$(1)$] If $\f: \X\ra \Y$ is a morphism in $\Kb{A}$, then, for each $i\in \Z$, $F(\f)=\chi_{i,F}(\f)+\g$ with $\g\in \Zmap_{i+1}(A)(\X, \Y)$. In particular, $F(\f)=\f+\h$ with $\h\in \sum_{j\in \Z}\Zmap_j(A)(\X, \Y)$.
\item[$(2)$] If $\cpx{f}:\X\ra\Y$ satisfies that $f^i\neq 0$ for at most one integer $i$, then $F(\f)=\f$.

\item[$(3)$] Let $n\geq 0$. Assume that $F$ fixes $\K^{\ell\leq n}(A\proj)$. Then
\begin{enumerate}
\item[{\rm (a)}] $F$ fixes $\K^{\ell\leq n+1}(A\proj)$ if and only if $F$ fixes $\K^{0,n+1}(A\proj)$.
\item[{\rm (b)}] For every morphism $\f: \X\ra \Y$ between objects of $\K^{0,n+1}(A\proj)$, there is a $\g\in \bigcap_{1\leq j\leq n+1}\Zmap_j(A)(\X, \Y)$ such that $F(\f)=\f+\g$.
\end{enumerate}
\end{enumerate}
\end{lemma}
\begin{proof}
(1) Applying $F$ to diagram $(*)$, the first statement follows from diagram $(**)$ and Lemma \ref{st}(4). For the second, choose $m\leq n$ such that both $X^i$ and $Y^i$ are zero whenever $i$ is outside $[m,n]$. Then $F(\f)=\chi_{m,F}(\f)+\g_{m+1}$ with $\g_{m+1}\in \Zmap_{m+1}(A)(\X, \Y)$. Since $F$ is a pseudo-identity, $F(f^{\leq m})=f^{\leq m}$, that is, $\chi_{m,F}(\f)^{\leq m}$ is given by $f^m$. By induction, we can assume that $F(f^{\geq m+1})=f^{\geq m+1}+\cpx{h'}$ for some $\cpx{h'}\in \sum_{i\in\mathbb{Z}}\Zmap_i(A)(X^{\geq m+1},Y^{\geq m+1})$. It is easy to see that each morphism in $\Zmap_i(A)(X^{\geq m+1},Y^{\geq m+1})$ can be extended to a morphism in $\Zmap_i(A)(\X,\Y)$ by adding a zero map from $X^m\ra Y^m$. Recall that $\chi_{m,F}(\f)^{\geq m+1}=F(f^{\geq m+1})$. It follows that $\chi_{m,F}(\f)=\f+\cpx{\tilde{h'}}$, where $\cpx{\tilde{h'}}$ is the extended morphism from $\cpx{h'}$ by adding a zero map from $X^m\ra Y^m$. Thus $F(\f)=\f+\cpx{g}+\cpx{\tilde{h'}}$ and $\cpx{g}+\cpx{\tilde{h'}}\in\sum_{i\in\mathbb{Z}}\Zmap_i(A)(\X,\Y)$.

(2) Assume $i\in\mathbb{Z}$ such that $f^j=0$ for all $j\neq i$, and consider  the following complex
$$\U:\;\;\cdots\lra X^{i-1}\lraf{d_X^{i-1}}X^i\lraf{f^i} Y^i\lraf{d_Y^i}Y^{i+1}\lra \cdots$$
with $U^i=X^i$.
We have the morphism $\uu{d_U^i}: U^{\leq i}[-1]=X^{\leq i}[-1]\ra U^{\geq i+1}=Y^{\geq i}[-1]$. Then $\f=\theta(\uu{d_Y^{i-1}})(\uu{d_U^i}[1])\pi(\uu{d_X^i})$. Since $F$ fixes all truncated triangles, $F$ fixes $\theta(\uu{d_Y^{i-1}})$, $\uu{d_U^i}[1]$ and $\pi(\uu{d_X^i})$. So $F$ fixes $\f$.

(3) (a) The necessity is clear. Now assume that $F$ fixes $\K^{0,n+1}(A\proj)$. Since the connecting isomorphism is $\id: F[1]\ra [1]F$, $F$ fixes $\K^{0,n+1}(A\proj)[i]$ for all $i\in\Z$. Now let $\f:\X\ra\Y$ be a nonzero morphism in $\K^{\ell\leq n+1}(A\proj)$ such that not both $\X$ and $\Y$ belong to $\K^{0,n+1}(A\proj)[i]$ for the same $i$. Let $s$ (respectively, $r$) be the minimal (respectively, maximal) degree  such that both $\X$ and $\Y$ have nonzero term. Then $r-s\leq n$.  By (1), $F(\f)=\chi_{s-1,F}(\f)+\g_s$ with $\g_s\in \Zmap_s(A)(\X, \Y)=0$. It follows that $F(\f)=\chi_{s-1,F}(\f)$. Note that $\chi_{s-1,F}(\f)^{\geq s}=F(f^{\geq s})$ and $F(f^{\geq s})=\chi_{r,F}(f^{\geq s})+\g_{r+1}$ with $\g_{r+1}\in \Zmap_{r+1}(A)(\X, \Y)=0$. However, \[\chi_{r,F}(f^{\geq s})^{\leq r}=F\left((f^{\geq s})^{\leq r}\right)=F(f^{[s,r]})=f^{[s,r]}\]
since $f^{[s, r]}\in \K^{\ell\leq n}(A\proj)$. Moreover, $\chi_{r,F}(f^{\geq s})^{\geq r+1}=F((f^{\geq s})^{\geq r+1})=0$. This implies that
\[f^{\geq s}=\chi_{r,F}(f^{\geq s})=F(f^{\geq s})=\chi_{s-1,F}(\f)^{\geq s}.\]
Clearly $\chi_{s-1,F}(\f)^{\leq s-1}=F(f^{\leq s-1})=0=f^{\leq s-1}$. It follows that $F(\f)=\chi_{s-1,F}(\f)=\f$.

(b) Let $0\leq j\leq n$. Observe that $f^{\leq j}$ and $f^{\geq j+1}$ belong to $\K^{\ell\leq n}(A\proj)$. Hence $F$ fixes $f^{\leq j}$ and $f^{\geq j+1}$, so $\chi_{j,F}(\f)=\f$. The assertion now follows from (1).
\end{proof}

With the help of Theorem \ref{t0} and Lemma \ref{coro}, we can give short proofs of some known results on Rickard's question.

\begin{exam}\label{ex}
(1) Rickard's question holds for hereditary algebras \cite{MY}.

Indeed, suppose that $A$ is a hereditary algebra. Let $\X, \Y$ be two indecomposable complexes in $\K^{0,1}(A\proj)$. Then $\X=(X^0\raf{d_X^0}X^1)$ and $\Y=(Y^0\raf{d_Y^0}Y^1)$, where $d_X^0$ and $d_Y^0$ are injective. It is easy to see that $\Zmap_1(A)(\X, \Y)=0$. By Theorem \ref{t0} and Lemma \ref{coro}(3), Rickard's question holds for $A$.

(2) Let $\Omega$ be the set of triples $(r, n, m)$ of integers such that $1\le r\le n$ and $m\geq 0$. A finite-dimensional algebra $A$ is called a {\it derived-discrete algebra} \cite{BGS} if $A$ is derived equivalent to either a hereditary algebra of Dynkin type or a $k$-algebra $\Lambda(r, n, m)$ with $(r, n, m)\in \Omega$, where $\Lambda(r, n, m)$ is defined by the quiver
\[\xymatrix@R=0.4em{
&& &\mathop{\bullet}\limits_1\ar[r]^-{\alpha_1} &\mathop{\bullet}\limits_2\cdots \mathop{\bullet}\limits_{n-r-2}\ar[r]^-{\alpha_{n-r-2}} &\mathop{\bullet}\limits_{n-r-1}\ar[rd]^{\alpha_{n-r-1}} &\\
\mathop{\bullet}\limits_{-m}\ar[r]^-{\alpha_{-m}} &\mathop{\bullet}\limits_{-m+1}\cdots  \mathop{\bullet}\limits_{-1}\ar[r]^-{\alpha_{-1}}&\mathop{\bullet}\limits_{0}\ar[ru]^{\alpha_{0}}&&&&\mathop{\bullet}\limits_{n-r}\ar[ld]_-{\alpha_{n-r}}\\
 && &\mathop{\bullet}\limits_{n-1}\ar[lu]_-{\alpha_{n-1}} &\mathop{\bullet}\limits_{n-2}\ar[l]_-{\alpha_{n-2}}\cdots \mathop{\bullet}\limits_{n-r+2}&\mathop{\bullet}\limits_{n-r+1}\ar[l]_-{\alpha_{n-r+1}}&
 }\]
with relations $\alpha_0\alpha_{n-1}=\alpha_{n-1}\alpha_{n-2}= \cdots = \alpha_{n-r+1}\alpha_{n-r}=0$. It was proved, with the help of the description of morphisms of the derived categories of gentle algebras \cite[Theorem 7.4]{ALP}, that Rickard's question holds for derived-discrete algebras \cite{BC, CZ}.

For $A:=\Lambda(r, n, m)$ with $r\neq 1$, we give a short proof of this result. Actually, by \cite[Theorem 7.4]{ALP}, $\dim_k\Hom_{\Kb{A}}(\X, \Y)\leq 1$ for all indecomposable objects $\X$ and $\Y$ of $\Kb{A}$. Hence either $\Zmap_i(A)(\X, \Y)=0$ for all $i\in \Z$, or $\Hom_{\Kb{A}}(\X, \Y)=\Zmap_i(A)(\X, \Y)$ for an $i\in \Z$. By Theorem \ref{t0} and Lemma \ref{coro}(2)--(3), Rickard's question holds for $A$.
\end{exam}

\section{A series of examples of non-standard derived equivalences}\label{example}

In this section we prove Theorem \ref{t} by constructing infinitely many counterexamples to Rickard's question.

We first prove the following criterion for a derived auto-equivalence $(\id,\phi)$ to be a triangle functor.

\begin{prop}\label{p}
Let $A$ be an Artin algebra. Suppose that $\phi: \id [1]\ra [1]\id$ is a natural isomorphism on $\Ka$ such that $\phi_{X^{\geq n}}[-1]=(\phi_{\X}[-1])^{\geq n}$ for any object $\X\in \Ka$ and for any $n\in \mathbb{Z}$. Then the following are equivalent:
\begin{enumerate}
\item[$(1)$] $(\id, \phi): \Ka\ra \Ka$ is a triangle functor.
\item[$(2)$] For any morphism $\f: \X\ra \Y$ in $\Kb{A}$, there is a morphism $\h: \Y\ra \X$ in $\Kb{A}$
\[\begin{array}{c}\xymatrix{
\X\ar[r]^{\f}\ar[d]_{\phi_{\X}[-1]-1}&\Y\ar@{..>}[dl]_{\h} \\
\X\ar[r]^{\f}&\Y
}\end{array}\]
such that $\f\h\f=\f(\phi_{\X}[-1]-1)$ in $\Kb{A}$.
\end{enumerate}
\end{prop}
\begin{proof}
Consider the condition:

\medskip
$(a)$: For any morphism $f: \X\ra \Y$ in $\Ka$, there is a morphism $\h: \Y\ra \X$ in $\Ka$ such that $\f\h\f=\f(\phi_{\X}[-1]-1)$ in $\Ka$.

\medskip\noindent
We will prove $(1)\Leftrightarrow (a)\Leftrightarrow (2)$. For simplicity, we write $F$ for the pair $(\id, \phi)$.

$(1)\Rightarrow (a)$ Assume that $F$ is a triangle functor. Then, for every morphism $\f: \X\ra \Y$ in $\Ka$,  there is an isomorphism $\cpx{\psi}: \cone(\f)\ra \cone(\f)$ in $\Ka$ such that the following diagram
 \[\begin{array}{c}\xymatrix@R=1.5em{
F(\Delta(\f)):& \X\ar[d]^{1}\ar[r]^{\f}&\Y\ar[r]^-{\theta(\f)}\ar[d]^{1}& \cone(\f)\ar[r]^{\phi_{\X}\pi(\f)}
\ar@{..>}[d]^{\cpx{\psi}}&\X[1]\ar[d]^{1} \\
\Delta(\f):& \X\ar[r]^{\f}&\Y\ar[r]^-{\theta(\f)}& \cone(\f)\ar[r]^{\pi(\f)}&\X[1]
}\end{array}\]
commutes in $\Ka$.
By Lemma \ref{st}(3), there is a morphism $\h: \Y\ra \X$ in $\Ka$ such that $\f\h\f=\f(\phi_{\X}[-1])-\f=\f(\phi_{\X}[-1]-1)$ in $\Ka$.

$(a)\Rightarrow (1)$ Note that $\Ka$ is a Krull-Schmidt category. Every morphism in $\Ka$ can be decomposed into an isomorphism and a radical morphism. There are complexes $\X_1, \X_2, \Y_1, \Y_2$ and  isomorphisms $\cpx{s}_X: \X\ra \X_1\oplus \X_2$ and $\cpx{s}_Y: \Y\ra \Y_1\oplus \Y_2$ such that the following diagram commutes in $\Ka$:
 \[\begin{array}{c}\xymatrix{
\X\ar[d]_{\cpx{s}_X}\ar[r]^{\f}&\Y\ar[d]^{\cpx{s}_Y}\\
\X_1\oplus \X_2\ar[r]^{\left(\begin{smallmatrix}
\cpx{s}&0 \\
0&\cpx{r}
\end{smallmatrix}\right)}&\Y_1\oplus \Y_2,
}\end{array}\]
where $\cpx{s}: \X_1\ra \Y_1$ is an isomorphism and $\cpx{r}: \X_2\ra \Y_2$ is a radical morphism. Then $F(\Delta(\f))$ is a triangle if and only if $F\big(\Delta\left(\begin{smallmatrix}
\cpx{s}&0 \\
0&\cpx{r}
\end{smallmatrix}\right)\big)$ is a triangle. Clearly, $F(\Delta(\cpx{s}))$ is a triangle as $\cone(\cpx{s})$ is the zero object. To prove that $F(\Delta(\f))$ is a triangle, we only need to show that $F(\Delta(\cpx{r}))$ is a triangle.

Consider the morphism $\cpx{r}$. By the assumption, there is a morphism $\h: \Y_2\ra \X_2$ in $\Ka$ such that $\cpx{r}\h\cpx{r}= \cpx{r}(\phi_{\X_2}[-1]-1)$ in $\Ka$.
By Lemma \ref{st}(3)-(5), there is an isomorphism $\cpx{\psi}: \cone(\cpx{r})\ra \cone(\cpx{r})$ in $\Ka$ such that the following diagram
 \[\begin{array}{c}\xymatrix@R=1.5em{
F(\Delta(\cpx{r})):& \X_2\ar[d]^{1}\ar[r]^{\cpx{r}}&\Y_2\ar[r]^-{\theta(\cpx{r})}\ar[d]^{1}& \cone(\cpx{r})\ar[r]^{\phi_{\X_2}\pi(\cpx{r})}
\ar@{..>}[d]^{\cpx{\psi}}&\X_2[1]\ar[d]^{1} \\
\Delta(\cpx{r}):& \X_2\ar[r]^{\cpx{r}}&\Y_2\ar[r]^-{\theta(\cpx{r})}& \cone(\cpx{r})\ar[r]^{\pi(\cpx{r})}&\X_2[1]
}\end{array}\]
commutes in $\Ka$, this means that $F(\Delta(\cpx{r}))$ is a triangle.

$(a)\Rightarrow (2)$ This implication is clear.

$(2)\Rightarrow (a)$ Let $\f: \X\ra \Y$ in $\Ka$. Then there is an $n\in \mathbb{Z}$ such that $H^i(\X)=H^i(\Y)=0$ for all $i\leq n$. Consider the following two complexes
$$\X_0: \;\;\; \cdots \lra 0\lra M\lraf{u} X^n\lraf{d_X^n} X^{n+1}\lraf{d_X^{n+1}}X^{n+2}\lra \cdots$$
$$\Y_0:\;\;\; \cdots \lra 0\lra N\lraf{v} Y^n\lraf{d_Y^n} Y^{n+1}\lraf{d_Y^{n+1}} Y^{n+2}\lra \cdots$$
where $M\raf{u} X^n$ is the kernel of $X^n\raf{d_X^n} X^{n+1}$ and $N\raf{v} Y^n$ is the kernel of $Y^n\raf{d_Y^n} Y^{n+1}$.
Then
$$t_{\X, \Y}: \Hom_{\Ka}(\X, \Y)\lra \Hom_{\K^b(A\md)}(\X_0, \Y_0),\;\;\; \g\longmapsto (H^n(g^{\geq n}), g^{\geq n})$$
is an isomorphism  \cite[Lemma 2.2(1)]{HX10} where $H^n(g^{\geq n}): M=H^n(X^{\geq n})\ra H^n(Y^{\geq n})=N$ is the homology of $g^{\geq n}$. Moreover, it is easy to see that, if two morphisms $\g_0, \g_1: X^{\geq n}\ra Y^{\geq n}$ of complexes are homotopic, that is, $\g_0\sim \g_1$, then
$$(H^n(\g_0), \g_0)\sim (H^n(\g_1), \g_1): \X_0\lra \Y_0.$$
Consider $f^{\geq n}: X^{\geq n}\ra Y^{\geq n}$. By (2), there is a morphism $\h_0: Y^{\geq n}\ra X^{\geq n}$ in $\Kb{A}$ such that $f^{\geq n}\h_0 f^{\geq n}=(\phi_{X^{\geq n}}[-1]-1) f^{\geq n}$ in $\Kb{A}$. Then we have a complex morphism
$$\h_1=(H^n(\h_0), \h_0): \Y_0\lra \X_0.$$
Define $\h:=t^{-1}_{\Y, \X}(\h_1)\in \Hom_{\Ka}(\Y, \X)$.
Then
 \begin{equation*}
\begin{aligned}
 t_{\X, \Y}(\f\h\f)&=(H^n(f^{\geq n}\h_0 f^{\geq n}), f^{\geq n}\h_0 f^{\geq n})\\
 &=(H^n((\phi_{X^{\geq n}}[-1]-1)f^{\geq n}), (\phi_{X^{\geq n}}[-1]-1)f^{\geq n})=t_{\X, \Y}((\phi_{\X}[-1]-1)\f)
\end{aligned}
 \end{equation*}
in $\K^b(A\md)$, where the last equality follows from the condition $\phi_{X^{\geq n}}[-1]=(\phi_{\X}[-1])^{\geq n}$. So $\f\h\f=(\phi_{\X}[-1]-1)\f$ in $\Ka$.
\end{proof}

{\bf In the rest of this section}, let $k:=\mathbb{Z}/2\mathbb{Z}$ and $A$ a finite-dimensional local commutative Frobenius $k$-algebra. Denote by $\rad(A)$ and $\soc(A)$ the radical and socle of $A$, respectively. Assume that $0\neq s\in \soc(A)\subseteq \rad(A)$ and $a\in \rad(A)$. Thus we always have $(\soc(A))^2=0$.


Let $\{a_1, \cdots, a_n\}$ be a $k$-basis of $A$ with $a_n=s$. Let $\omega$ be the following $k$-linear map
$$\omega: A\lra k, \;\;\; \sum_{1\leq i\leq n}k_ia_i\longmapsto k_n.$$
Then $\omega$ is a Frobenius form of $A$.

\begin{lemma}\label{m} Keep the notations above. If $a$ satisfies $\omega(x^2)=\omega(a x)$ for all $x\in A$, then $\omega(\tr(Z^2))=\omega(\tr(a Z))$ for all square matrices $Z$ over $A$.
\end{lemma}
\begin{proof}
We calculate $\tr(Z^2)$ as follows:
$$\begin{aligned}
  \tr(Z^2) & =\sum_i(Z^2)_{i, i}=\sum_i\sum_jZ_{i, j}Z_{j, i}
  = \sum_{i<j} Z_{i, j}Z_{j, i} +\sum_{i>j} Z_{i, j}Z_{j, i}+\sum_i Z_{i, i}^2\\
  & = 2\sum_{i<j} Z_{i, j}Z_{j, i} +\sum_i Z_{i, i}^2 = \sum_{i} Z_{i,i}^2.
\end{aligned}$$
Hence
$\omega(\tr(Z^2))=\sum_i \omega(Z_{i, i}^2)=\sum_i \omega(a Z_{i, i})=\omega(\tr(a Z)).$
\end{proof}

\medskip
Now we define a natural transformation $\xi(a): \id[1]\ra [1]\id$ as follows.
For a complex $\Y$ in $\Ka$, we define $\xi(a)_{\Y}: \Y[1]\ra \Y[1]$ by
$$\xi(a)_{\Y}^i: Y^{i+1}\lra Y^{i+1},\;\; y\longmapsto (1+a)y.$$
Clearly, this is a morphism of complexes.

\begin{lemma}\label{k}
Assume $\omega(x^2)=\omega(a x)$ for all $x\in A$. Then, for each morphism $\f: \X\ra \Y$ in $\Kb{A}$, there is a morphism $\g: \Y\ra \X$ in $\Kb{A}$:
\[\begin{array}{c}\xymatrix{
\X\ar[r]^{\f}\ar[d]_{\xi(a)_{\X}[-1]-1}&\Y\ar@{..>}[dl]_{\g} \\
\X\ar[r]^{\f}&\Y
}\end{array}\]
such that $\f\g\f= \f(\xi(a)_{\X}[-1]-1)$.
\end{lemma}
\begin{proof}
Let $\f: \X\ra \Y$ be a morphism in $\Kb{A}$. We may assume that $\f$ is a representative of the corresponding homotopy class of $\f$, thus $\f$ is a morphism of complexes. We need to prove the following property:

\medskip
(P) There is a morphism $\g: \Y\ra \X$ of complexes such that $\f\g\f\sim \f(\xi(a)_{\X}[-1]-1)$.

\medskip
\noindent
By definition $\f(\xi(a)_{\X}[-1]-1)=a\f$, so property (P) is equivalent to saying that
$$a \f\in \{\f\g\f+d_D^{-1}\h\,|\,\g\in Z^0(\cpx{C}), \h\in D^{-1}\}=\f Z^0(\cpx{C})\f+B^0(\cpx{D}),$$
where $\cpx{C}:=\Hom^{\bullet}(\Y, \X)$ and $\cpx{D}:=\Hom^{\bullet}(\X, \Y)$.

We shall use the non-degenerate bilinear forms $\lan -, -\ran: C^0\times D^0\ra k$ and $\lan -, -\ran: D^0\times C^0\ra k$ defined in Lemma \ref{form}. Indeed, for any $\g\in C^0$ and $\cpx{z}\in D^0$, we have
$$\lan \f\g\f, \cpx{z}\ran =\sum_{i\in \mathbb{Z}}(-1)^i\omega\big(\tr(f^ig^if^iz^i)\big)=\sum_{i\in \mathbb{Z}}(-1)^i\omega\big(\tr(g^if^iz^if^i)\big)=\lan \g, \f\cpx{z}\f\ran.$$
Lemma \ref{form}(3) says that $Z^0(\cpx{C})^{\perp}=B^0(\cpx{D})$. Letting $\cpx{g}$ range over $Z^0(\cpx{C})$, we obtain
$$\f\cpx{z}\f\in B^0(\cpx{D})\iff \cpx{z}\in  (\f Z^0(\cpx{C})\f)^{\perp}.$$
We also have $B^0(\cpx{D})^{\perp}=Z^0(\cpx{C})$ by Lemma \ref{form} (3). We deduce from
\[\big(\f Z^0(\cpx{C})\f+B^0(\cpx{D})\big)^{\perp}=\big(\f Z^0(\cpx{C})\f\big)^{\perp}\cap \big(B^0(\cpx{D})\big)^{\perp}\]
that $\cpx{w}\in C^0$ belongs to $\big(\f Z^0(\cpx{C})\f+B^0(\cpx{D})\big)^{\perp}$ if and only if $\cpx{w}\in Z^0(\cpx{C})$ and $\f\cpx{w}\f\in B^0(\cpx{D})$. Now for such a $\cpx{w}$, one has
$$\lan a \f, \cpx{w}\ran =\sum_{i\in \mathbb{Z}}(-1)^i\omega(\tr(a f^iw^i))=\sum_{i\in \mathbb{Z}}(-1)^i\omega(\tr(f^iw^if^iw^i))=\lan \f\cpx{w}\f, \cpx{w}\ran=0,$$
where the second equality follows from Lemma \ref{m}, and the last one follows from Lemma \ref{form}(3). This implies that
$a\f\in \big(\f Z^0(\cpx{C})\f+B^0(\cpx{D})\big)^{\perp\perp}=\f Z^0(\cpx{C})\f+B^0(\cpx{D}).$
Thus $\f$ has property (P), and this finishes the proof.
\end{proof}

\begin{theorem}\label{3.1}
Suppose that $a^2\neq 0$ and $\omega(x^2)=\omega(a x)$ for all $x\in A$. Then
 $$(\id, \xi(a)): \Ka\lra \Ka$$
  is a non-standard derived equivalence.
\end{theorem}
\begin{proof}
For simplicity, we write $\xi$ for $\xi(a)$ and write $F:=(\id, \xi): \Ka\ra \Ka$. For $z\in A$, let $\rho(z)$ be the right multiplication map:
 $$\rho(z): A\lra A, x\longmapsto xz \mbox{ for } x\in A.$$

Clearly, for an object $\X$ in $\Ka$ and for $m\in \mathbb{Z}$, the equality $\xi_{X^{\geq m}}[-1]=(\xi_{\X}[-1])^{\geq m}$ holds. By Proposition \ref{p} and Lemma \ref{k}, $F$ is a derived equivalence. Assume to the contrary that $F$ is standard. Since $F$ is a pseudo-identity, it must be naturally isomorphic to  $(\id, \id)$ as triangle functors. Let $\mu: F=(\id,\xi)\ra (\id, \id)$ be a natural isomorphism of triangle functors. We have a commutative diagram in $\Ka$:
  \[\begin{array}{c}\xymatrix@R=1.5em{
F(A[1])=A[1]\ar[r]^{\xi_{A}}\ar[d]_{\mu_{A[1]}} &F(A)[1]=A[1]\ar[d]^{\mu_{A}[1]}\\
\id(A[1])=A[1]\ar[r]^{\id}&\id(A)[1]=A[1].
}\end{array}\]
Then $\mu_{A[1]}=\mu_A[1]\circ \xi_{A}$. We may assume $\mu_A:=\rho(u)$ for a unit $u\in A$. Then $\mu_{A[1]}=\rho(u(1+a))[1]$.

Let $C(a):=\cone(\rho(a))$ and $C(s):=\cone(\rho(s))$. Since $\mu$ is a natural isomorphism, we have a commutative diagram in $\Ka$:
\[\xymatrix@R=1.5em{
\Delta(\rho(a)): &A\ar[r]^{\rho(a)}&A\ar[d]^{\mu_{A}}\ar[r]^-{\theta(\rho(a))} &C(a)\ar[d]^{\mu_{C(a)}}\ar[r]^{\pi(\rho(a))}&A[1]\ar[d]^{\mu_{A[1]}} \\
\Delta(\rho(a)): &A\ar[r]^{\rho(a)}&A\ar[r]^-{\theta(\rho(a))} &C(a)\ar[r]^{\pi(\rho(a))}&A[1].}\]
By Lemma \ref{st}(4), $\mu_{C(a)}=(\rho(u), \rho(u+a x))$ for some $x\in A$ with $a xa=0$. It follows from $a^2\neq 0$ that $x\in \rad(A)$. Similarly, we have $\mu_{C(s)}=(\rho(u(1+a)), \rho(u+\lambda s))$ with $\lambda\in k$.

Since $a^2\neq 0$,  the cyclic $A$-module $Aa$ is nonzero. Thus the socle of $Aa$ is $\soc(A)$, and therefore there is $y\in A$ such that $y a=s$. Since $a^2\neq 0$ and $(\soc(A))^2=0$, we know $a\notin\soc(A)$ and $y\in \rad(A)$. Consider the morphism
 $$\phi:=(\rho(1), \rho(y)): C(a)\lra C(s)$$
in $\Ka$. As $\mu$ is a natural isomorphism, we have the commutative diagram in $\Ka$:
  \[\xymatrix@R=1.5em{
C(a)\ar[rr]^{\phi}\ar[d]_{\mu_{C(a)}}&&C(s)\ar[d]^{\mu_{C(s)}}\\
C(a)\ar[rr]^{\phi} &&C(s).}\]

We prove that this diagram is actually not commutative, and hence a contradiction appears. Indeed, since both $y$ and $x$ are in $\rad(A)$, $\lambda s y=0$ and $a x y= x s=0$. Then
$$\mu_{C(s)}\phi-\phi\mu_{C(a)}=(\rho(u(1+a)), \rho(u y+\lambda s y))-(\rho(u), \rho(u  y+a x y))=(\rho(ua), 0)$$ in $\Ka$.
However, since $u$ is a unit in $A$, the morphism
$$(\rho(u a), 0): C(a)\lra C(s)$$
of complexes is never zero-homotopic, and consequently $\mu_{C(s)}\phi\neq \phi\mu_{C(a)}$ in $\Ka$, that is, the above-mentioned diagram is not commutative.
\end{proof}

{\it Remark}. (1) The restricted triangle equivalence $(\id, \xi(a)): \Kb{A}\ra \Kb{A}$ is also non-standard. Indeed, the proof of Theorem \ref{3.1} shows that the obstruction to standard equivalence already appears in $\Kb{A}$.

(2) Let us explain why we take $k=\mathbb{Z}/2\mathbb{Z}$. The key point is the  condition that  $\omega(x^2)=\omega(ax)$ holds for all $x\in A$. This happens only if $k=\mathbb{Z}/2\mathbb{Z}$. Actually, if $x\in A$  and $\omega(x^2)=\omega(ax)\neq 0$, then, for each $\lambda\in k$, we get $\lambda^2\omega(x^2)=\omega((\lambda x)^2)=\omega(a\lambda x)=\lambda\omega(ax)$. It follows that $\lambda^2=\lambda$ for all $\lambda\in k$. This can happen only in $\mathbb{Z}/2\mathbb{Z}$.

(3) Since we take $k=\mathbb{Z}/2\mathbb{Z}$, both sides of  the equation $\omega(x^2)=\omega(a x)$ are $k$-linear in $x$. This means that we only need to check the equation for $x$ running over a $k$-basis of $A$.


\medskip
Let $A, B$ be  finite-dimensional local commutative Frobenius $k$-algebras with the Frobenius forms $\omega_A$ and $\omega_B$, respectively. Assume that $k$ is a splitting field for both $A$ and $B$.  Then the tensor product algebra $A\otimes_k B$ is still a finite-dimensional local commutative Frobenius $k$-algebra with the Frobenius form defined by
$$\omega_A\otimes \omega_B: A\otimes_k B\lra k,\;\; \sum_ix_i\otimes y_i\longmapsto \sum_i\omega_A(x_i)\omega_B(y_i).$$

\begin{prop} \label{prop4.6} Let $A,B$ be as above, and let $a\in A$ and $b\in B$ such that $a^2\neq 0$ and $\omega_A(x^2)=\omega_A(ax)$ for all $x\in A$; $b^2\neq 0$ and $\omega_B(y^2)=\omega_B(b y)$ for all $y\in B$. Let $\Lambda:= A\otimes_k B$ be the tensor product of $A$ and $B$ over $k$.
Then
$$(\id, \xi(a\otimes b)): \KL\lra \KL$$
 is a non-standard derived equivalence.
\end{prop}
\begin{proof}
Clearly, $(a\otimes b)^2\neq 0$. Consider elements of $A\otimes_k B$ in the form $x\otimes y$ with $x\in A$ and $y\in B$. For $x\otimes y\in A\otimes_kB$, we have
$$\begin{aligned}
  (\omega_A\otimes \omega_B)((x\otimes y)^2) & =(\omega_A\otimes \omega_B)(x^2\otimes y^2)=\omega_A(x^2)\omega_B(y^2)\\
  & =\omega_A(a x)\omega_B(b y)=(\omega_A\otimes \omega_B)((a\otimes b)(x\otimes y))
\end{aligned}$$
Since $k=\Z/2\Z$, the map $z\mapsto z^2$ is $k$-linear. It follows that $(\omega_A\otimes \omega_B)(z^2)= (\omega_A\otimes \omega_B)((a\otimes b)z)$ for all $z\in A\otimes_kB$. By Theorem \ref{3.1}, we have the assertion.
\end{proof}

Proposition \ref{prop4.6} shows that once we know a counterexample to Rickard's question, we can construct infinitely many counterexamples by tensor products of algebras.

\medskip
The following is a concrete example satisfying conditions in Theorem \ref{3.1}. Thus
Theorem \ref{t} becomes a consequence of the following example and Proposition \ref{prop4.6}.

\begin{exam}\label{counterexa}
Let $m\geq 1$ and $n_i\geq 1$ for all $1\leq i\leq m$.  Let
$$A:=k[x_1,x_2,\ldots,x_m]/(x_1^{2n_1+1},x_2^{2n_2+1},\ldots,x_m^{2n_m+1}).$$
Then $A$ is the tensor product of $A_i:=k[x_i]/(x_i^{2n_i+1})$ for $1\le i\le m$. For each $i$, let $a_i=x_i^{n_i}$ and let $\omega_i:A_i\ra k$ be the linear map extracting the coefficient of $x_i^{2n_i}$. Then, for each $x_i^t$, it is easy to check that
$\omega_i((x_i^t)^2)=\omega_i(x_i^{n_i+t})=\omega_i(x_i^{n_i}x_i^t).$
It follows from Theorem \ref{3.1} that $(\id,\xi(a_i)):\K^{-,b}(A_i\proj)\ra \K^{-,b}(A_i\proj)$ is a non-standard derived equivalence for all $i$. Taking tensor products, we deduce by Proposition \ref{prop4.6} that $(\id,\xi(a_1a_2\cdots a_m)):\Ka\ra\Ka$ is a non-standard derived equivalence.
\end{exam}

We end this section by the following conjecture.

\smallskip
\begin{conject} Rickard's question holds for finite-dimensional algebras over fields of characteristic $p\ne 2$.
\end{conject}

 \smallskip
{\bf Acknowledgements:} The research work was supported partially by the National Natural Science Foundation of China (Grant 12501044) and Beijing Natural Science Foundation (No. 1252011). The third named author is supported by the China Postdoctoral Science Foundation (Grant 2025M783067). During a visit of CCXi to Southern University of Science and Technology in August 2026, a preliminary version of the work was completed. Xi is very grateful to Professor Zhicheng Feng for his invitation and warm hospitality.

When considering the proof of Lemma \ref{k} in the case $A=k[x]/(x^3)$, we asked ChatGPT Pro for help. It suggested a similar condition.

\medskip
\begin{thebibliography}{abcd}
\bibitem{ALP} K. K. Arnesen, R. Laking, D. Pauksztello, Morphisms between indecomposable complexes in the bounded derived category of a gentle algebra. J. Algebra {\bf 467} (2016) 1-46.

\bibitem{BGS} G. Bobi\'{n}ski, C. Gei\ss, A. Skowro\'{n}ski, Classification of discrete derived categories. Cent. Eur. J. Math. {\bf 2} (2004) 19-49.

\bibitem{BC} G. Bobi\'nski, T. Ciborski, Derived equivalences for the derived discrete algebras are standard. Algebr. Represent. Theory {\bf 29} (2026) 331-352.

\bibitem{C} X. W. Chen, A note on standard equivalences. Bull. Lond. Math. Soc. {\bf 48} (2016) 797-801.

\bibitem{C1} X. W. Chen, Representability and autoequivalence groups. Math. Proc. Cambridge Philos. Soc. {\bf 171} (2021) 657-668.

\bibitem{CY} X. W. Chen, Y. Ye, The {$\mathbf{D}$}-standard and {$\mathbf{K}$}-standard categories. Adv. Math. {\bf 333} (2018) 159-193.

\bibitem{CZ} X. W. Chen, C. Zhang, The derived-discrete algebras and standard equivalences. J. Algebra {\bf 525} (2019) 259-283.

\bibitem{H} Happel, D.: \emph{Triangulated categories in the representation theory of finite-dimensional algebras}. London Mathematical Society Lecture Note Series, 119. Cambridge University Press, Cambridge, 1988.

\bibitem{HX10} W. Hu, C. C. Xi, Derived equivalences and stable equivalences of Morita type. I. Nagoya Math. J. {\bf 200} (2010) 107-152.

\bibitem{Krause} H. Krause, Krull--Schmidt categories and projective covers,
Expos. Math. \textbf{33} (2015), no.~4, 535-549.

\bibitem{MY} J. I. Miyachi, A. Yekutieli, Derived Picard groups of finite-dimensional hereditary algebras. Compositio Math. {\bf 129} (2001) 341-368.

\bibitem{N} A. Neeman, Some new axioms for triangulated categories. J. Algebra {\bf 139} (1991) 221-255.

\bibitem{R0} J. Rickard, Morita theory for derived categories. J. Lond. Math. Soc. {\bf 39} (1989) 436-456.

\bibitem{R} J. Rickard, Derived equivalences as derived functors. J. Lond. Math. Soc. {\bf 43} (1991) 37-48.

\bibitem{RZ} R. Rouquier, A. Zimmermann, Picard groups for derived module categories. Proc. Lond. Math. Soc. {\bf 87} (2003) 197-225.

\bibitem{Y}  A. Yekutieli, Dualizing complexes, Morita equivalences and the derived Picard group of a ring. J. London Math. Soc. {\bf 60} (1999) 723-746.
\end{thebibliography}

\end{document}
